% =====================================================================
%  Capitolo 18 — Il metodo persona e i player type: Maslow, Bartle,
%  Keirsey, Bateman, Caillois, Lazzaro e il modello prismatico
%  (Game Design, primo anno)
%  Fonte: trascrizione della lezione GD18 (Game Design, primo anno),
%  stesso docente dei capitoli 9, 10 e 15 (corso di design methods), di
%  cui la lezione è la diretta prosecuzione: si chiude la lezione sulle
%  decisioni del capitolo 15 e si apre il tema dell'audience. Slide
%  della lezione: DM06 "Persona & Types" (slide di Diego Ricchiuti,
%  2026). Lezione in aula, con alcuni studenti collegati su Zoom.
%  Contenuti, nell'ordine della lezione: l'apertura (telecamere accese,
%  niente più intrusioni durante la lezione, comunicare è il mantra del
%  game design, il gioco con Keanu Reeves e Dr. Langeskov); la piramide
%  dall'agency alla mastery (feedback a breve termine, conseguenze a
%  lungo termine, Dark Souls e i Gargoyle); il gioco si fa per qualcun
%  altro; come capire l'audience: statistica, metodo persona, player
%  type; data informed e non data driven, lo skinning e lo unique
%  selling point; la digressione sulla piramide di Maslow, i suoi
%  limiti (le notti in bianco, il treno, i bisogni "vitaminici"), cibo,
%  sesso e violenza (Dead Space, Arancia meccanica, The Witcher, Cooking
%  Mama, gli spot dei fast food), i tre gradini che i videogiochi
%  toccano; le motivazioni di Quantic Foundry e le differenze culturali
%  (Cina e Stati Uniti); il metodo persona in breve; la tassonomia di
%  Bartle (la favola della buonanotte, i MUD, gli assi, i colori, i semi
%  delle carte, le quattro definizioni; la domanda sui cheater e
%  l'amico degli achievement di Mirror's Edge; intento esplicito e
%  implicito; single player e multiplayer); la domanda sull'uso della
%  tassonomia; la pausa; i temperamenti di Keirsey e le sedici
%  personalità; Bartle-Keirsey; il modello demografico di Bateman e il
%  casual di oggi; perché allargare Bartle e i suoi limiti; Caillois
%  (agon, alea, mimicry, ilinx, il quiz sull'abbinamento, il D&D del
%  docente, l'alea che rende tutti pari, il poker); Lazzaro e le quattro
%  chiavi del divertimento; la tabella del modello prismatico; GNS+ e
%  MDA+ da non usare; Handy, Gallup, Covey, motivazione, problem
%  solving e obiettivo finale; il "non task": rifare il modello unificato
%  (le case di Hogwarts di Federica, le classi e le personificazioni);
%  l'analisi per generi (la leggenda della tripla A, Binding of Isaac,
%  Pong, Ubisoft, Slay the Spire); le feature di gameplay per player
%  type (la sorpresa come moltiplicatore, Minecraft, le emote di Clash
%  Royale e Marvel Snap); i segmenti di Quantic Foundry (Celeste,
%  Farewell, Metroid Dread, Spelunky, Animal Crossing e Tom Nook, Just
%  Cause, Clash of Clans); i tre usi del modello (analisi, suggerimenti
%  di design con le biciclette, soluzione di un problema: il gioco di
%  carte e il bluff, con l'interruzione di un collega); la pausa; il
%  persona finale; la vision rivista e i design value; la profilazione
%  del target audience; il calendario e la chiusura.
%  Integrazioni da "The Art of Game Design" (Jesse Schell) e "21st
%  Century Game Design" (Chris Bateman e Richard Boon), segnalate come
%  approfondimenti e tenute distinte dalla lezione.
%
%  Pacchetti richiesti nel preambolo del libro (già presenti nel
%  main.tex): \usepackage{booktabs}, \usepackage{amsmath},
%  \usepackage{enumitem}, \usepackage{tikz} con
%    \usetikzlibrary{shapes.geometric,shapes.symbols,shapes.misc,
%      arrows.meta,positioning,calc,fit}
%  (le librerie vanno caricate nel preambolo, non nel capitolo).
%  Tutte le etichette di questo capitolo hanno il prefisso "gd18:"
%  (gd18:sec:..., gd18:fig:..., gd18:tab:...) e tutti gli stili TikZ
%  il prefisso "gd18", per non collidere con gli altri capitoli.
%
%  Nota sulla fonte: la trascrizione automatica è buona nelle parti
%  discorsive, ma storpia moltissimi nomi propri e termini inglesi:
%  "Bartol", "Barto", "Battle", "Bartolomeo" per Bartle; "Chirsey",
%  "Kirsey", "Kirsten" per Keirsey; "Calwa", "Callo A", "Caloua" per
%  Caillois; "Lazzar" per Lazzaro; "Civer", "Cibere", "Asciugher",
%  "Shiver", "Chiver" per achiever; "Mini craio", "Mimicai" per
%  mimicry; "Inx", "E-Lynx", "Unix" per ilinx; "Unix endpoint" per
%  unique selling point; "Udema di Koster" per l'atomo di Koster;
%  "Chianu Yis", "Chiamurides" per Keanu Reeves; "Autoful Boyfriend"
%  per Hatoful Boyfriend; "Tom Dacity Life" per Tomodachi Life; "Zed
%  Spire" per Slay the Spire; "Metroid Red" per Metroid Dread; "Alflife"
%  per Half-Life; "BanduScape" probabilmente per RuneScape; "Airstone"
%  per Hearthstone; "Gacia" per gacha. Dove il riferimento è
%  ricostruibile con buona sicurezza lo si riporta nella forma corretta
%  (segnalandolo in nota quando non è ovvio); dove non lo è, lo si
%  segnala in nota invece di indovinare. Come nei capitoli 15 e 17,
%  l'attesa iniziale e le due pause sono state riempite dal
%  riconoscimento vocale con frasi in loop o captate da una fonte
%  esterna ("Grazie a tutti per l'attenzione che ci ha spiato",
%  "Qual è la risposta che ci può essere?", frasi sul "governo di Dalai
%  Lama" e sulla "migliore dal 2005"): sono artefatti, omessi; lo stesso
%  vale per un "il cento e il cento..." ripetuto decine di volte durante
%  una domanda. Le figure sono le slide originali della lezione; le
%  tabelle ne trascrivono il contenuto testuale.
% =====================================================================

\chapter[Persona e player type]{Il metodo persona e i player type: da
Maslow e Bartle al modello prismatico}
\label{gd18:cap:player-type}

Il capitolo~\ref{gd15:cap:decisioni} si era chiuso sulle decisioni nel gioco e su un
task singolare: un gioco comico, una ``truffa'' che sovverte le aspettative. Questa
lezione, tenuta dallo stesso docente del corso di \emph{design methods}, chiude
quel discorso e ne apre uno nuovo, che il docente presenta senza giri di parole come
``forse una delle cose più importanti che vedremo e che vedrete quest'anno'': i
\textbf{player type}. Le slide della lezione si intitolano \emph{Persona \& Types}, e
il titolo dice già le due anime della lezione: il \textbf{metodo persona}, cioè la
costruzione di un individuo fittizio che riassume il target audience, e i
\textbf{tipi di giocatore}, cioè un modo di approssimare e categorizzare persone
molto diverse a partire da come giocano e da che cosa cercano nel gioco.

Il filo conduttore è sempre lo stesso, ed era già comparso nelle lezioni sulla
ricerca (capitolo~\ref{cap:ricerca-idee}) e sulla vision
(capitolo~\ref{gd9:cap:agency-visione}): \textbf{il gioco si fa per qualcun altro}.
Per farlo bene bisogna capire chi è quel qualcun altro, che cosa cerca e perché; e
siccome i dati che si raccolgono su persone diverse sono ``potenzialmente molto
lontani l'uno dall'altro, difficili da correlare'', serve un modo per
\textbf{approssimarli}. La lezione costruisce questo modo un pezzo alla volta: parte
dalla tassonomia di \textbf{Richard Bartle}, che va saputa ``a memoria'', e la allarga
con altri modelli (\textbf{Keirsey}, \textbf{Bateman}, \textbf{Caillois},
\textbf{Lazzaro} e altri ancora) fino a un \textbf{modello unificato}, che il
coordinatore del corso chiama \textbf{modello prismatico}: tante facce diverse dello
stesso prisma, che guardano sempre la stessa cosa. Poi mostra come usarlo nel
lavoro di tutti i giorni, in tre casi concreti.

La struttura della lezione, e quindi del capitolo, è questa:
\begin{enumerate}
  \item l'apertura e la chiusura della lezione sulle decisioni
        (sezione~\ref{gd18:sec:apertura});
  \item la \textbf{piramide} che va dall'agency alla mastery, e il gioco fatto per
        qualcun altro (sezione~\ref{gd18:sec:piramide});
  \item i tre strumenti per capire l'audience, e il principio \emph{data informed}
        (sezione~\ref{gd18:sec:strumenti});
  \item la digressione sulla \textbf{piramide di Maslow}, sui suoi limiti e su cibo,
        sesso e violenza (sezione~\ref{gd18:sec:maslow});
  \item le motivazioni di Quantic Foundry e le differenze culturali
        (sezione~\ref{gd18:sec:quantic});
  \item il \textbf{metodo persona}, in breve (sezione~\ref{gd18:sec:persona});
  \item la \textbf{tassonomia di Bartle}, nelle sue tre letture
        (sezione~\ref{gd18:sec:bartle});
  \item i \textbf{temperamenti di Keirsey}, le sedici personalità e la coppia
        Bartle-Keirsey (sezione~\ref{gd18:sec:keirsey});
  \item il \textbf{modello demografico di Bateman} e i tipi intermedi
        (sezione~\ref{gd18:sec:bateman});
  \item il \textbf{modello prismatico}: Caillois, Lazzaro, le colonne da non usare e
        le altre (sezione~\ref{gd18:sec:prismatico});
  \item il ``non task'': rifarsi il modello (sezione~\ref{gd18:sec:rifare});
  \item l'\textbf{analisi per generi} e le \textbf{feature di gameplay} per player
        type (sezioni~\ref{gd18:sec:generi} e~\ref{gd18:sec:feature});
  \item i \textbf{segmenti di Quantic Foundry} (sezione~\ref{gd18:sec:segmenti});
  \item i \textbf{tre usi} del modello: analisi, suggerimenti di design, soluzione
        di un problema (sezione~\ref{gd18:sec:casi});
  \item il \textbf{persona finale} (sezione~\ref{gd18:sec:persona-finale});
  \item la \textbf{vision} rivista e i \emph{design value}
        (sezione~\ref{gd18:sec:vision});
  \item la \textbf{profilazione del target audience}, il calendario e la chiusura
        (sezioni~\ref{gd18:sec:ta} e~\ref{gd18:sec:chiusura}).
\end{enumerate}

% ---------------------------------------------------------------------
\section{L'apertura: comunicare, e chiudere le decisioni}
\label{gd18:sec:apertura}

\subsection{Telecamere accese, e niente più intrusioni}

La lezione si tiene in aula, con qualcuno collegato da remoto. Si comincia con la
prova dell'audio: il docente chiede a Federica di testarlo (``sento tutto forte e
chiaro''), e gli altri ``eventualmente arriveranno e noi saremo pronti ad
accoglierli''. Poi ripete una richiesta già fatta: \textbf{telecamere accese}. Uno
studente avverte che sta finendo di mangiare; nessun problema, ``ho visto
letteralmente persone mangiare'', e si è liberi di farlo (magari, scherza il docente,
senza mangiare ``il pezzo di pasta davanti'' alla camera).

Segue un discorso che il docente aveva già fatto prima della lezione e il venerdì
precedente. Ha parlato con Alessandro, e ci sono stati alcuni cambiamenti, ``primo
fra tutti il fatto che non ci siano più \textbf{invasioni} durante la vostra ora di
lezione'': un risultato pensato per il benessere degli studenti, la direzione verso
cui lui e gli altri docenti si muovono sempre, cioè dare ``l'ecosistema e l'anno
educativo migliore possibile''. Per questo ripete l'invito: se ci sono dubbi,
difficoltà o problemi, anche individuali, bisogna parlarne, tra studenti o con
qualcuno che li riporti; le questioni più delicate si affrontano dopo la lezione.
Quel giorno il docente si ferma apposta in aula ``fino a quando volete'', ma in
realtà è disponibile ogni giorno in cui è presente, faccia a faccia. Il motivo è un
principio che vale ben oltre l'aula:
\begin{quote}
Non si possono affrontare dei problemi se non si comunica. È questo poi il mantra
del game design.
\end{quote}

\subsection{Il gioco con Keanu Reeves, e un consiglio}

Non ci sono domande sui task (che il docente non ha ancora commentato: ``lo farò al
più presto''). Resta il ``grandissimo, incredibile'' gioco con \textbf{Keanu Reeves}
del task precedente (capitolo~\ref{gd15:sec:task}), che il docente spera tutti
abbiano potuto ``godere''. A chi l'esperienza è piaciuta, ripete il consiglio: un
gioco più serio dello stesso creatore, \emph{Dr.\ Langeskov, The Tiger, and The
Terribly Cursed Emerald},\footnote{Nella trascrizione: ``Dr. Langslob, la tigre,
l'Emerald Impazzito''. È lo stesso titolo già consigliato nel
capitolo~\ref{gd15:sec:task}.} gratuito e molto più breve di \emph{The Stanley
Parable} (che è ``blasonato'' e che tutti conoscono). Un modo per ``espandere
l'esperienza'', anche se ``ahimè, senza Keanu Reeves come protagonista principale''.

Il docente chiede se ci sono dubbi sulle decisioni, sui tipi di decisione, sulle
\emph{golden rule}; poi, visto che nel frattempo è arrivata Francesca (e anche
Francesco, salutato con un ``nessun problema''), ripete la domanda. Nessun
dubbio: si \textbf{saltano le golden rule} e si conclude la lezione sulle decisioni.

% ---------------------------------------------------------------------
\section{Dall'agency alla mastery, e per chi si fa il gioco}
\label{gd18:sec:piramide}

\subsection{La piramide}

Le ultime due lezioni hanno parlato di \textbf{agency} e di \textbf{decisioni}, che
stanno alla base di una piramide che il docente riprende a voce, gradino per
gradino.\footnote{La piramide non è tra le slide di questa lezione: il docente la
mostra dalle slide della lezione precedente (DM05, capitolo~\ref{gd15:cap:decisioni}).}
\begin{description}
  \item[Agency.] Il gradino più basso: i \textbf{limiti} entro i quali il giocatore
        decide. ``Non parliamo più di libertà'': come sempre, il docente invita a
        usare il linguaggio tecnico il più possibile (la libertà non esiste, esiste
        l'agency, capitolo~\ref{gd9:cap:agency-visione}).
  \item[Decisioni.] Qualsiasi cosa il giocatore scelga o non scelga deve avere un
        \textbf{outcome}, e qui il docente rimanda all'atomo di
        Koster\footnote{Nella trascrizione: ``il famoso Udema di Koster''. Si tratta
        quasi certamente dell'atomo di gioco di Koster
        (capitolo~\ref{cap:problema-soluzione}), in cui il feedback chiude il ciclo
        tra azione del giocatore e modello mentale.} e a tutto ciò che si è detto sui
        feedback: l'outcome va \textbf{segnalato}, deve essere chiaro, deve essere
        comunicato.
  \item[Feedback e conseguenze a breve termine.] Un gradino sopra: ``ho compiuto la
        scelta, mi viene detto che la scelta è stata fatta''. Il famoso banner del
        tipo ``si ricorderà della tua scelta''\footnote{Nella trascrizione: ``lei si
        ricorderà, Madeline, si ricorderà della scelta''. È con ogni probabilità il
        ``Clementine will remember that'' dei giochi Telltale già discusso nel
        capitolo~\ref{gd15:cap:decisioni}; il nome è storpiato.}, quel prompt di UI,
        ricade qui, nelle \emph{short term consequences}. E vale a qualsiasi livello:
        si spara con un'arma e c'è il \emph{rumble}; si prende danno e il personaggio
        lampeggia, il cuoricino si rompe.
  \item[Conseguenze a lungo termine.] L'aumento o la diminuzione di statistiche e le
        meccaniche che ne derivano stanno più sopra. Come il feedback viene
        \emph{capito} dal giocatore, del resto, ``sta un gradino sopra''.
  \item[Mastery.] In cima, l'obiettivo finale: acquisire la \textbf{padronanza dei
        sistemi}. È un discorso complesso, che verrà ripreso con un altro framework
        (``di cui non è che sono proprio grande fan, però questo è un buon punto'').
        Non si parla qui di padronanza di una singola meccanica o di una combo, ma
        di \textbf{imparare le regole del gioco}.
\end{description}

\subsection{La mastery forma il giocatore}

Sulla mastery il docente si sofferma. Il designer crea il gioco per un target
audience specifico, ma è giusto dire anche il contrario: \textbf{attraverso
l'esperienza che il designer crea, il giocatore evolve}. Alla fine dell'esperienza
il giocatore diventa, idealmente, ciò che il designer pensa che possa essere.

L'esempio è \emph{Dark Souls}. Quando i boss ``vi massacrano'', quando si incontra
il Demone Capra o il Demone Toro\footnote{Nella trascrizione: ``il demone di Thor o il
demone capra''; sono il \emph{Taurus Demon} e il \emph{Capra Demon} del primo
\emph{Dark Souls}.} o, ``quello più infame'', i \textbf{Gargoyle}, e si muore un
certo numero di volte, si impara a capire il pattern dei boss, a navigare
l'ambiente, a gestire le fiaschette di Estus. Questo è già \textbf{spingere il
modello mentale} del giocatore nella direzione di ciò che il designer ritiene il
giocatore ideale di \emph{Dark Souls}. Non deve diventare per forza il pro player
definitivo; ma ``tramite come lo educate'' lo si sta spingendo in una direzione e si
sta formando la sua mente in quella direzione. La mastery in cima alla piramide
parla di questo, ad alto livello. Ci sono poi forme di mastery più basse (capire
quando schivare, come gestire la stamina in un combattimento, riconoscere
l'\emph{anticipation} di un attacco nemico mentre lo carica per evitarlo), ma
servono solo a spiegare la piramide. Con quanto visto finora, si è arrivati al
gradino del feedback a breve termine.

\subsection{Il gioco si fa per qualcun altro}

Come già detto più volte, tutti ``stiamo facendo un gioco per qualcun altro''.
Potrebbe essere per sé stessi, ma allora c'è il \textbf{rischio d'impresa}: il gioco
potrebbe non vendere. Chi è assunto in un'azienda, generalmente, non fa il gioco per
sé, a meno che l'azienda non sia sua e ne sia ``forse l'unico dipendente''. Si fa il
gioco per un audience di persone, e le persone sono diverse: chi gioca ai party
game sforzandosi ``incredibilmente per premere il tasto A'' sul telecomando della
Wii; il ragazzino tedesco che spacca il monitor perché ha perso a
\emph{Counter-Strike}; il pro player pieno di coppe e medaglie; il divorziato che
gioca tutto il giorno a \emph{Candy Crush}. Qualsiasi demografica si scelga,
qualsiasi target audience, bisogna avere chiaro \textbf{verso quale TA si sta
facendo il gioco}, in quale direzione si stanno spingendo le risorse, e accertarsi
sempre che quella direzione sia corretta. Un primo passo in questo senso si è già
fatto con la vision, con i documenti di visione e con la ricerca. Ora si fa il passo
successivo.

% ---------------------------------------------------------------------
\section{Come capire l'audience: tre strumenti}
\label{gd18:sec:strumenti}

\subsection{Statistica, persona, player type}

Come capire l'audience? Una parte la si è già vista: tutto ciò che riguarda la
\textbf{statistica}, i dati effettivi di ricerca. Oggi se ne vedono altre due:
\begin{enumerate}
  \item il \textbf{metodo persona}, importante ``relativamente'': un altro modo di
        far confluire i dati, concretizzandoli in un'ipotetica persona. Il docente
        chiede se qualcuno ne ha già sentito parlare; poi spiega che si tratta di
        creare un individuo \textbf{fittizio}, il cui profilo raccoglie tutti i dati
        trovati nella fase di ricerca, in modo da avere un individuo ideale verso
        cui si fa il gioco. Questa scheda si stampa, si appende, ``ce l'abbiamo
        sempre chiaro esattamente come le vision'': quando c'è un problema, quando
        si discute, si fa riferimento a questo individuo che non esiste. Di solito
        se ne fa più di uno, e lo fanno soprattutto le aziende medio-grandi. È un
        modo per avere sempre chiara la direzione e soprattutto \emph{per chi} si fa
        il gioco, ``senza poi chiamare in causa la personalità'';
  \item il \textbf{metodo dei player type}, il vero argomento del giorno: ``una cosa
        molto fica. Non voglio creare troppo hype, però è molto fico questo metodo'',
        che aiuta enormemente non solo il processo di design ma molti processi della
        creazione di un gioco.
\end{enumerate}

La statistica gli studenti l'hanno già vista: raccogliere dati e come farlo
diventerà ancora più importante quando dovranno creare un gioco e fare davvero la
ricerca, trovare i dati e scegliere. Sempre nella direzione della vision, cercando di
concretizzare i dati in metodi diversi.

\subsection{Data informed, non data driven}

Il docente richiama la parte finale della lezione sui feedback
(capitolo~\ref{gd10:cap:feedback}): a meno di lavorare per King o per qualche gioco
mobile specifico, \textbf{nel 90\% dei casi} il lavoro del designer è \textbf{data
informed}, non \textbf{data driven}. Ciò significa che \textbf{non} si prendono i dati
per copiare ciò che è già stato fatto, né si copia esattamente ciò che l'utente dice
di volere per inserirlo nel gioco. Lo fanno pochi studi estremamente grossi, con una
mole enorme di dati, tanto da poter fare \textbf{clustering} (come visto nella
lezione sui flowchart e sui grafici, capitolo~\ref{cap:brainstorming-flowchart});
servono risorse enormi nella ricerca e negli strumenti, che già un team medio-grande
non ha: ``doppia A? No, no, no''.

Ciò che interessa è fare ricerca su ciò che l'utente vuole e cercare di capire,
\emph{al di là di ciò che chiede}, da cosa è stimolato e che cosa cerca, per trovare
la \textbf{soluzione creativa migliore} a un problema individuato: quello che gli
studenti hanno già cominciato a fare con il task sulla ricerca.

\paragraph{Lo skinning.} A volte, comunque, il design si copia, e non è quello il
problema. Si fa \textbf{skinning} di un gioco: di solito non per intero (se lo si fa
per intero è una scelta aziendale, generalmente di risparmio), ma su alcune
meccaniche specifiche, perché costa di meno o perché lo standard dell'industria è
quello e lo si adotta. Ma lo \textbf{unique selling point} deve stare altrove: ``se
copiamo lo unique selling point il nostro gioco crolla'', il prodotto tendenzialmente
non vende, indipendentemente dalla grandezza dell'azienda. ``Vedete Ubisoft nei tempi
recenti.''

Chiusa questa parte, il programma del giorno è chiaro: il metodo persona, ``una
piccola parentesi'', e soprattutto i player type, ``perché sono la cosa più
importante''.

\slidefig{img/GD/gd18-maslow}{La piramide dei bisogni di Maslow: i videogiochi toccano i tre gradini più alti.}{slide \emph{DM06 -- Persona \& Types}}{gd18:fig:maslow}

% ---------------------------------------------------------------------
\section{Una digressione: la piramide di Maslow}
\label{gd18:sec:maslow}

\subsection{I cinque gradini}

La \textbf{piramide di Maslow} (figura~\ref{gd18:fig:maslow}) è ``un metodo ormai molto
blasonato'' per comunicare la \textbf{gerarchia dei bisogni umani}: più in alto si
va, più il bisogno è motivante e trainante per la persona, ma i bisogni dei gradini
precedenti devono essere già risolti. Il docente la percorre velocemente dal basso:
\begin{enumerate}
  \item \textbf{bisogni fisiologici}: respirare, bere, mangiare, avere un tetto sulla
        testa, dormire (e dormire al sicuro), avere dei vestiti, il sesso e la
        riproduzione;
  \item \textbf{sicurezza}: sicurezza personale, avere introiti e quindi soldi da
        spendere, un impiego, risorse in generale, la salute, la proprietà privata;
  \item \textbf{amore e appartenenza}: amici, gruppi sociali, famiglia, intimità con
        un partner, senso di connessione con altre persone;
  \item \textbf{stima} (\emph{esteem}): essere riconosciuti dagli altri per il proprio
        status, per ciò che si è e si fa; rispetto, forza, libertà, ``insomma stiamo
        ampliando lo spettro'';
  \item \textbf{autorealizzazione} (\emph{self-actualization}): cercare di diventare
        la versione migliore possibile di sé. Non c'è un parametro: si tratta di
        andare, nella propria vita, nella direzione che si desidera. ``Sarebbe bello
        pensare che il fatto che voi siate qui in questo momento è per questa
        direzione qui'': diventare la migliore versione di sé facendo il gioco dei
        propri sogni. ``Io ancora non ci sono arrivato, quindi la strada è lunga.
        Ognuno ha i propri metri di misura.''
\end{enumerate}

\subsection{Non è l'unico modello, e non spiega tutto}

Prima di arrivare al punto della digressione, il docente premette che la piramide di
Maslow \textbf{non è l'unico modo} di guardare ai bisogni umani. È consolidata perché
ha una lunga storia di condivisione ed è probabilmente il primo risultato che si
trova cercando ``bisogni umani''; ma non spiega tutto.

Per dimostrarlo chiede in aula: quanti non hanno mai dormito una notte? Diverse mani
si alzano, e il docente chiede il perché. Capita anche di non mangiare: lui stesso
quel giorno ha pensato di saltare il pranzo per non perdere il treno (``per poco non
ho perso il treno. Non è che giochi molto a mio favore questa storia, però è andata
così''). Se in un certo momento si ritiene più importante un bisogno superiore e si
considerano risolti quelli inferiori, allora il voler restare in contatto con altre
persone, il voler giocare con altri o il voler giocare \textbf{al day one} a un gioco
può \textbf{sovrascrivere} il bisogno di dormire o di mangiare. La teoria, per quanto
condivisa, ha delle falle, ``come molte cose d'altronde''.

Ci sono altri modelli. Il docente ne cita uno: i cosiddetti \textbf{bisogni
``vitaminici''}, che vedono ciascuno di questi bisogni come una necessità
vitaminica. Si va a dormire quando non ce la si fa più, o quando si ritiene di aver
consumato abbastanza energie; non tutti mangiano a pranzo e a cena, c'è chi per
lavoro o per altre necessità ha orari ``sballati''. Il punto è che \textbf{quando si
ritiene che uno di questi bisogni si stia riducendo, ci si muove per risolverlo}.

\subsection{Cibo, sesso e violenza}

C'è poi un'altra osservazione, che tornerà. È vero che più in alto si va più i
bisogni sono motivanti, ma ``rimaniamo delle bestie con i peli'': ci sono tre cose
che non tanto muovono l'essere umano, quanto ne \textbf{catturano l'attenzione}:
\textbf{il cibo, il sesso e la violenza}. Sesso e violenza, nota il docente, sono
alla base di due delle industrie fondanti dell'umanità: la guerra e l'industria del
sesso.

\paragraph{Dead Space.} Quanti hanno giocato a \emph{Dead Space}? Qual è la cosa più
memorabile? Gli studenti citano i necromorfi e i doppiaggi (``certamente i doppiaggi
sono incredibili''); il docente cerca una scena che viene citata sempre, anche da chi
non sa che viene da \emph{Dead Space}: la \textbf{scena dell'occhio} e
dell'ago.\footnote{La scena in cui il protagonista deve infilarsi un ago nell'occhio
è in \emph{Dead Space 2} (2011). Nella trascrizione la discussione è confusa e non si
capisce se il docente la attribuisca al primo o al secondo capitolo.} È
particolarmente memorabile perché si attacca a questo senso di \textbf{pericolo e di
violenza}. È lo stesso motivo per cui esistono \emph{Postal}, \emph{Il Grezzo 2} ``o
altri giochi che non citerò''; ed è il meccanismo su cui funziona \emph{Arancia
meccanica}. Il docente mostra la scena avvisando chi è sensibile (``adesso si inizia
a vedere l'ago'').

\paragraph{Sesso e cibo.} Lo stesso vale per le scene di sesso in \emph{The Witcher}.
Per il cibo, \emph{Cooking Mama} non è un grande esempio (``non è che sono
particolarmente invitanti''), ma le grandi cucine di \emph{Final Fantasy} sì; e
funziona anche nella vita di tutti i giorni: le pubblicità dei fast food con
l'enorme panino.

\paragraph{Trappole luccicanti, con un limite.} Tutte queste cose catturano
l'attenzione: sono ``trappole molto luccicanti'', ma sono facili e \textbf{non si
possono usare sempre}. La violenza non si può usare costantemente, altrimenti si
riduce il \textbf{tasso di assorbimento}: una volta che attirano l'attenzione e la
persona se ne rende conto, la loro efficacia si riduce notevolmente. È un tema che
verrà ripreso al secondo anno.\footnote{Durante questa parte uno studente propone un
esempio legato all'universo SCP (un metodo di contenimento particolarmente violento,
forse da \emph{SCP -- Containment Breach}); la trascrizione è troppo frammentaria per
ricostruirlo, e il docente stesso non conosce il caso.}

\subsection{I tre gradini che i videogiochi toccano}

Il punto della digressione: \textbf{i videogiochi possono toccare i primi tre
gradini dall'alto} (figura~\ref{gd18:fig:maslow}).
\begin{itemize}
  \item \textbf{Autorealizzazione}: c'è chi la trova nei videogiochi tramite
        l'\textbf{escapismo}, e non solo: c'è chi fa del giocare, in maniera
        competitiva, la propria carriera.
  \item \textbf{Amore e appartenenza}: il bisogno di relazionarsi con altri può
        essere soddisfatto collaborando con persone vere o con NPC. Senza entrare in
        discorsi complicati, c'è chi trova un sufficiente senso di connessione con un
        NPC: si pensi alle \emph{romance}, o a \emph{Hatoful Boyfriend}, già visto a
        lezione, il cui obiettivo è allacciare una relazione con un piccione; per
        qualcuno può risolvere davvero un bisogno di appartenenza.
  \item \textbf{Stima}: riconoscimento, status, rispetto.
\end{itemize}
Il discorso, precisa il docente, vale \emph{adesso}: nessuno dice che fra qualche
anno giocando non ci si possa ``riempire la pancia''. È già borderline: praticando
gli esport ad alti livelli si può avere un impiego; ma lì si esce dal gioco, pur
restando dentro l'esperienza. Infine, mettersi nei panni di un altro per provare
emozioni, sentirsi un pilota di Formula 1 o un cavalcatore di draghi, è ciò che si
è visto come \textbf{esperienza vicaria}. Con la premessa che è un discorso
circoscritto a questo modello, i videogiochi toccano e risolvono questi tre gradini.

\slidefig{img/GD/gd18-motivazioni}{Le motivazioni di gioco di Quantic Foundry, divise in sei gruppi.}{slide \emph{DM06 -- Persona \& Types}}{gd18:fig:motivazioni}

% ---------------------------------------------------------------------
\section{Le motivazioni di Quantic Foundry, riprese}
\label{gd18:sec:quantic}

Altri modelli categorizzano le necessità e le motivazioni umane in gruppi diversi.
Uno è quello di \textbf{Quantic Foundry}, già incontrato nel
capitolo~\ref{cap:ricerca-idee}, di cui il docente consiglia ``caldamente'' il talk
alla GDC (a dicembre ne consiglierà altri). Le motivazioni sono raccolte in sei
gruppi (figura~\ref{gd18:fig:motivazioni}): \textbf{action}, \textbf{social},
\textbf{mastery}, \textbf{achievement}, \textbf{immersion} e \textbf{creativity}.
Sono cose già viste, che vanno ricordate, ``perché poi questo inizia a diventare
parte di un calderone molto importante''.

Le slide mostrano anche le differenze di genere (il docente invita a ``lasciar
perdere'' il genere) e le tre tabelle di \emph{gameplay analysis} già viste, con le
preferenze agli estremi di ogni motivazione. Il punto su cui il docente insiste è
un altro: le motivazioni sono uno \textbf{spettro}, e variano molto con la
\textbf{demografica} e la \textbf{cultura}. Una persona europea, nord-europea, avrà
tendenzialmente motivazioni più sviluppate di un'altra rispetto a un asiatico o a un
americano. Il grafico Stati Uniti contro Cina (figura~\ref{gd18:fig:usa-cina})
mostra per esempio che la Cina ricerca di più la \textbf{challenge}, mentre gli Stati
Uniti hanno un percentile più alto nel \textbf{completamento}.

\slidefig{img/GD/gd18-usa-cina}{Le motivazioni di gioco a confronto: Stati Uniti e Cina.}{slide \emph{DM06 -- Persona \& Types}}{gd18:fig:usa-cina}

E qui sta il problema. Persone diverse hanno necessità diverse e sono spinte da
motivazioni diverse; finora si è visto come fare ricerca e ottenere dati, ma sono
dati ``potenzialmente molto lontani l'uno dall'altro, difficili da correlare''. Se
l'obiettivo è raggiungere un individuo ideale verso cui si crea il gioco, serve
\textbf{un certo tipo di approssimazione}. Ed è questa approssimazione l'argomento
della lezione.

% ---------------------------------------------------------------------
\section{Il metodo persona, in breve}
\label{gd18:sec:persona}

\slidefig{img/GD/gd18-persona}{Il metodo persona: un utente fittizio basato sul target audience.}{slide \emph{DM06 -- Persona \& Types}}{gd18:fig:persona}

Il \textbf{persona method} funziona molto bene e non si usa solo nell'industria dei
videogiochi. Il docente lo tratterà meglio alla fine della lezione
(sezione~\ref{gd18:sec:persona-finale}); per ora ne dà ``solo gli estremi''
(figura~\ref{gd18:fig:persona}).

\begin{definition}[Persona]
Un \textbf{persona} è il profilo di un utente \textbf{fittizio}, costruito a partire
dalla ricerca di pre-produzione sul target audience, verso cui si fa il gioco. Non è
basato su una persona reale, \textbf{né sul designer}: è un'amalgama di tutto ciò
che si sa del proprio target audience.
\end{definition}

Si parte dalla ricerca già fatta, tutta quella di pre-produzione, e si cerca di
creare un vero e proprio ``pedigree'': il profilo di una persona ipotetica. Per
farlo bisogna rispondere a ``un macello di domande'', tutte quelle già viste nella
lezione sulla pre-produzione (i \textbf{nove perché}, capitolo~\ref{cap:ricerca-idee}):
che età ha, che genere, che livello di cultura, qual è la sua routine, che tipo di
famiglia, quanto tempo ha a disposizione, quanto gioca e quanto no, che cosa gli
piace e che cosa no. Le slide insistono:
\begin{quote}
\emph{Details are EVERYTHING in the P.\ method}: dire che a lui o a lei piace
studiare è generico. Che cosa gli piace studiare? Quando? Perché?
\end{quote}

Tutto questo deve essere basato sul target audience: non è una personalità e non è
basato su un design, ``sempre per il discorso che voi fate un gioco per qualcun
altro''. Una volta costruito il persona, \textbf{tutte le decisioni} di gioco
rispetto alla vision dipendono da lui: da ciò che gli può piacere (\emph{likes}) e
da ciò che non gli può piacere (\emph{dislikes}). Nel profilo devono esserci tutti i
suoi \textbf{bisogni} e i suoi \textbf{problemi}, tutto ciò che definisce
concretamente una persona: leggendone la scheda, lo si deve avere in testa.

% ---------------------------------------------------------------------
\section{La tassonomia di Bartle}
\label{gd18:sec:bartle}

\subsection{La favola della buonanotte dei designer}

Ciò che interessa arrivare a vedere il più velocemente possibile è il \textbf{metodo
dei player type}: approssimare i dati e \textbf{categorizzare archetipi di
giocatori}. Il docente lo introduce come ``la storia della buonanotte dei
designer''.

Tanto tempo fa si cominciò a pensare a come approssimare gruppi di persone molto
diverse in base a \emph{come} giocano e a \emph{che cosa cercano} nel gioco. Il
``buon avventuriero'' \textbf{Richard Bartle} esaminò i \textbf{multi-user dungeon},
i cosiddetti \textbf{MUD}: giochi ormai scomparsi, antenati di quelli che oggi
sarebbero i giochi con enormi quantità di utenti attivi.\footnote{Il docente dice
``quello che in tempi moderni sarebbero i MOBA''. I MUD sono giochi di ruolo
testuali multigiocatore in tempo reale; i loro eredi diretti sono più propriamente i
MMORPG.} A partire da ciò che i giocatori facevano nel gioco, e compilando un
profilo degli utenti con un \textbf{survey}, una serie di domande, Bartle divise i
giocatori in \textbf{quattro categorie}. Sono queste le quattro categorie da cui
parte tutto il resto della lezione, per arrivare a qualcosa di molto più grande: un
\textbf{modello di approssimazione dell'utenza} che possa guardare gli stessi tipi
di giocatori da più punti di vista. È ciò che si fa con i framework: unire più modi
di vedere gli stessi giocatori. Diego\footnote{Diego Ricchiuti, autore delle slide
del corso di \emph{design methods}.} lo chiama \textbf{modello prismatico};
precedentemente veniva chiamato \textbf{modello unificato}: diverse facce di un
prisma, che guardano sempre la stessa cosa.

Il docente mostra il \textbf{paper} di Bartle e lo condivide su Zoom (chiedendo a
Federica di girarlo agli altri), insieme al \textbf{test} per i più curiosi:
``chiaramente è molto brutto da vedersi''. La slide, invece, ``è già bella,
vestita''.\footnote{Il paper è Richard Bartle, \emph{Hearts, Clubs, Diamonds,
Spades: Players Who Suit MUDs}, 1996. Il test online che ne deriva è il cosiddetto
\emph{Bartle Test of Gamer Psychology}.}

\slidefig{img/GD/gd18-bartle}{La tassonomia di Bartle: gli assi agire/interagire e giocatori/mondo, e i quattro tipi.}{slide \emph{DM06 -- Persona \& Types}}{gd18:fig:bartle}

\subsection{Gli assi e i quattro tipi}

Bartle disegna il suo modello su due assi (figura~\ref{gd18:fig:bartle}):
\begin{itemize}
  \item un asse va da \textbf{acting} (agire) a \textbf{interacting} (interagire).
        ``A mio modesto avviso, opinione personalissima'', la cosa importante non è
        l'\emph{interacting} in sé, ma l'\textbf{interacting with}, e dall'altra
        parte non l'\emph{acting} ma l'\textbf{acting on}: già la preposizione dice
        qualcosa del giocatore;
  \item l'altro asse va dai \textbf{players} (gli altri giocatori) al
        \textbf{world} (il mondo di gioco).
\end{itemize}
Dalle risposte alle domande di Bartle emergono quattro tipi di giocatori:
\begin{itemize}
  \item i \textbf{killer}: \emph{acting on players};
  \item gli \textbf{achiever}: \emph{acting on world};
  \item gli \textbf{explorer}: \emph{interacting with world};
  \item i \textbf{socializer}: \emph{interacting with players}.
\end{itemize}

\paragraph{I colori.} La slide aggiunge un passo successivo ``per far quadrare i
conti'': un colore per tipo. I \textbf{killer sono rossi}, gli \textbf{achiever
blu}, gli \textbf{explorer verdi}, i \textbf{socializer gialli}. Bartle non li aveva
categorizzati così, ma con i \textbf{semi delle carte} (``picche, quadri, fiori e
l'altro, che non ricordo'');\footnote{I quattro semi del titolo del paper: gli
achiever sono i quadri ($\diamondsuit$), gli explorer le picche ($\spadesuit$), i
socializer i cuori ($\heartsuit$) e i killer i fiori ($\clubsuit$). Si veda anche
Jesse Schell, \emph{The Art of Game Design}, capitolo 9.} siccome è potenzialmente
confusionario, il docente consiglia di restare legati alla colorazione delle slide e
di usarla in futuro anche per eventuali \emph{color coding} legati ai player type,
nella documentazione e non solo.

\paragraph{L'iconografia.} Non è l'unico modo di rappresentarli. Cercando
``Bartle taxonomy'' si trovano illustrazioni ricorrenti: i killer con le pistole, i
socializer ``bonaccioni'' che parlano e comunicano, l'explorer come grande
esploratore con la lente d'ingrandimento, l'achiever con la sacca di soldi o i
trofei.

\subsection{Chi sono}

\paragraph{Killer: agire sui giocatori.} L'obiettivo dei killer è avere
\textbf{potere} e usarlo contro gli altri: impattare gli altri, modificare il modo
in cui agiscono nel gioco, \textbf{interferire} con il funzionamento del mondo di
gioco o con l'esperienza degli altri giocatori. All'estremo, vogliono ``mettere i
bastoni tra le ruote'' agli altri. Questo comporta una serie di emozioni e di
necessità specifiche, ma basta sapere che l'obiettivo ultimo è \textbf{agire sugli
altri giocatori}. E lo fanno \textbf{a prescindere dalle regole}: a ogni costo,
anche uscendo dal gioco (non nel senso di ``bruciare la casa alle persone'', ma
usando \emph{cheat} o metodi paralleli).

\paragraph{Achiever: agire sul mondo.} L'achiever agisce sul mondo
\textbf{secondo le regole} del mondo di gioco: resta dentro i suoi confini, e quei
limiti, nella sua testa, lo tutelano. Il suo obiettivo è \textbf{accumulare
sicurezza} e \textbf{status}: accumulare tutto l'ottenibile nel mondo di gioco. Qui
ricadono i \textbf{completisti}, il cui obiettivo è ottenere tutti i trofei (``quella
è la sicurezza''), e chi fa \textbf{grinding}: livellare, ottenere parti di mostro
rare ripetendo continuamente la stessa missione.

\paragraph{Explorer: interagire con il mondo.} Molto simile, ma con una differenza
chiave: se il guadagno principale dell'achiever è fisico, dentro il mondo di gioco,
quello dell'explorer è \textbf{mentale}. L'achiever guadagna oggetti, l'explorer
guadagna \textbf{informazioni}. Vuole sapere \textbf{come funziona} il mondo di
gioco: che cosa porta un bioma, che cosa porta il bioma innevato, che cosa si può
fare combinando oggetti o abilità, se una torcia accesa sotto un acquazzone o sotto
una cascata si spegne o no.

\paragraph{Socializer: interagire con i giocatori.} L'obiettivo del socializer è
\textbf{giocare insieme}: il gioco di gruppo, fare party, comunicare con altre
persone, potenzialmente anche fare \emph{role playing}; in generale conoscere altre
persone, o conoscere la storia di altri NPC.

\begin{table}[ht]
\centering
\small
\begin{tabular}{@{}p{0.14\textwidth}p{0.2\textwidth}p{0.08\textwidth}p{0.48\textwidth}@{}}
\toprule
\textbf{Tipo} & \textbf{Asse} & \textbf{Colore} & \textbf{Definizione (slide)} \\
\midrule
Killer & acting on players & rosso & interferiscono con il funzionamento del mondo
di gioco o con l'esperienza degli altri giocatori \\
Achiever & acting on world & blu & accumulano status battendo le sfide del mondo di
gioco basate sulle regole \\
Explorer & interacting with world & verde & scoprono i sistemi che fanno
funzionare il mondo di gioco \\
Socializer & interacting with players & giallo & stringono relazioni con altri
giocatori o con NPC \\
\bottomrule
\end{tabular}
\caption{I quattro tipi di Bartle.}
\label{gd18:tab:bartle}
\end{table}

\subsection{Una domanda: i cheater sono killer?}

Uno studente chiede se chi usa i cheat in \emph{GTA} ricada nei killer, e perché lo
faccia chi bara in un gioco dove poi gli altri ``muoiono''.\footnote{La domanda è
trascritta in modo molto confuso; il senso ricostruibile è questo.} ``Dipende, perché
questa è una bella domanda'', e porta a un punto fondamentale:
\begin{quote}
Nessuno è veramente una sola di queste cose, e spesso noi cambiamo il player type a
seconda del gioco di riferimento.
\end{quote}
La domanda da farsi è \textbf{perché il giocatore usa quel mezzo}. Il motivo per cui
si è parlato di motivazioni è proprio questo: bisogna capire che cosa cerca il
giocatore.

\paragraph{L'amico degli achievement.} Il docente ha un amico ossessionato dagli
achievement e dal battere ogni record (``sono sicuro che è entrambe queste due
cose''). In \emph{Mirror's Edge} alcuni achievement non sono più ottenibili, perché
hanno chiuso i server; lui se li è sbloccati con i cheat in dieci minuti. Perché?
``Per dormire tranquillo la notte.'' L'azione in sé (barare) è più legata al killer;
ma \textbf{l'azione, la meccanica, ci dice poco dei giocatori}: ciò che ci dice
qualcosa è \textbf{come la uso e perché}. Se l'obiettivo è ottenere tutti gli
achievement, si è probabilmente nel regno dell'achiever.

\paragraph{I cheat dell'explorer.} Allo stesso modo, c'è chi usa i cheat (per
esempio in \emph{Monster Hunter} o in \emph{Dark Souls})\footnote{Nella trascrizione:
``Mostrante''; la lettura \emph{Monster Hunter} è probabile ma non certa.} per capire
dove si appiattisce la \textbf{curva delle statistiche}: per fare il
\textbf{min-maxing} di cui si è parlato la volta scorsa, cioè capire quando diventa
meno efficiente investire punti in forza, visto che l'aumento di una statistica è
proporzionalmente sempre meno vantaggioso. Oppure per guardare le animazioni e
vedere esattamente in quale frame c'è l'invincibilità, o per studiare le interazioni
dei \emph{Sims}. Usare strumenti alternativi per ottenere informazioni può ricadere
nel campo dell'explorer. ``Nessuno è solamente socializer, nessuno è solamente
achiever.''

\slidefig{img/GD/gd18-bartle-intenti}{Le sottocategorie di Bartle: intento esplicito e implicito.}{slide \emph{DM06 -- Persona \& Types}}{gd18:fig:bartle-intenti}

\subsection{Prima lettura: intento esplicito e implicito}

La tassonomia si può guardare da diversi punti di vista. Il primo è quello
dell'\textbf{intento esplicito o implicito} (figura~\ref{gd18:fig:bartle-intenti}),
che dà a ogni tipo due sottotipi con nomi propri.\footnote{È la suddivisione che
Bartle stesso propone nel suo modello a otto tipi, in \emph{Designing Virtual Worlds}
(2003).} Di nomi, avverte il docente, se ne vedranno ``davvero tanti''. Nel lavoro
con il team, di solito, per semplicità si usano i \textbf{nomi classici}: dire
``killer'' significa usare il nome di Bartle, ma portandosi dietro tutto ciò che si
dice da qui in avanti (politician e griefer ricadono comunque lì dentro). Se poi
interessa un frammento specifico (per esempio la grossa reputazione del politician),
si cita il sottotipo. Quando si categorizza un gioco, si individuano i player type
principali o si discute tra designer se una meccanica piaccia al target audience, ci
si appoggia ai nomi classici. ``Tuttavia, non dimenticate tutti gli altri'': è il
motivo stesso di questa lezione.

\begin{description}
  \item[Killer.] \textbf{Politician} (esplicito): vuole una grossa
        \emph{buona} reputazione, a tutti i costi. \textbf{Griefer} (implicito): il
        suo vago obiettivo è una grossa \emph{cattiva} reputazione. Sono l'uno
        l'opposto dell'altro, ma con lo stesso risultato: essere riconosciuti, essere
        visti, avere un impatto sugli altri. Il griefer è ``quello che trolla'': non
        vuole vincere né ottenere il risultato che il designer vuole; vuole vedere il
        compagno piangere, la persona soffrire, gli anni di lavoro di un giocatore di
        \emph{Minecraft} andare in fumo. Vuole quel potere, applicato sugli altri.
  \item[Achiever.] \textbf{Planner} (esplicito): si pone un obiettivo e punta a
        raggiungerlo; ``lista di obiettivi e si va''. Una meccanica che si sposa bene
        con l'achiever è un \textbf{set di goal chiaro}: il completista, quando deve
        completare un gioco, guarda la lista degli achievement, che cosa c'è da fare e
        come si fa. \textbf{Opportunist} (implicito): si guarda intorno in cerca di
        cose da fare, ma non sa quali siano finché non le trova.
  \item[Explorer.] \textbf{Scientist} (esplicito): è metodico nell'acquisire
        conoscenza, ha una lista di metodi, ``esattamente come funziona il metodo
        scientifico'': testa una serie di cose finché le sue ipotesi non sono
        validate o confutate. \textbf{Hacker} (implicito): ha una comprensione
        \textbf{intuitiva} del mondo virtuale, senza bisogno di testare le sue idee:
        in un sandbox prova una soluzione perché ``sicuramente acqua e fuoco non vanno
        bene insieme'', o perché ``magma e acqua creano un nuovo tipo di materiale e
        ci posso costruire''.
  \item[Socializer.] \textbf{Networker} (esplicito): cerca e \emph{valuta} le
        persone con cui vale la pena stare, i gruppi di amici con cui stare insieme.
        \textbf{Friend} (implicito): gode semplicemente della compagnia degli altri.
\end{description}

\slidefig{img/GD/gd18-bartle-sp-mp}{I tipi di Bartle in single player e in multiplayer.}{slide \emph{DM06 -- Persona \& Types}}{gd18:fig:bartle-sp-mp}

\subsection{Seconda lettura: single player e multiplayer}

Il secondo punto di vista è \textbf{single player} contro \textbf{multiplayer}
(figura~\ref{gd18:fig:bartle-sp-mp}).

\paragraph{Killer.} In single player l'obiettivo è \textbf{dominare l'intelligenza
artificiale} e il mondo di gioco. Qui si vede che, a seconda di ciò che si osserva,
killer e achiever possono avere lo stesso obiettivo: essere \textbf{primi nella
leaderboard} è ``comune e goloso'' per entrambi. La differenza è \emph{come} lo si
ottiene: il killer fa qualsiasi cosa, anche barare, anche agire disonestamente;
l'achiever è agonistico, sviluppa la mastery dei sistemi e diventa il migliore
perché ``ha sudato, perché si è allenato'', e il numero uno è la conseguenza dei suoi
sforzi. Il multiplayer è invece ``il sale della vita del killer'': dominare le IA può
essere semplice, ma la sfida più grande, perché imprevedibile, sono i
\textbf{giocatori}. Misurare le proprie abilità contro un altro giocatore, il
\emph{thrill of the hunt}, cacciarli, tracciarli, l'impulso di imporsi. È
esattamente ciò che offrono \emph{Dead by Daylight}, \emph{Evolve} e in parte
\emph{Among Us}: il killer vuole essere il \emph{big bad boy}, l'aggressore,
l'assassino, perché ha più potere. È una situazione impari in cui l'assassino ha più
strumenti e può letteralmente togliere dalla partita gli altri giocatori.

\paragraph{Achiever.} In single player: ottenere tutto ciò che si può ottenere, il
\textbf{100\%} del gioco; in generale avere un obiettivo e risolverlo metodicamente,
per esempio con il grinding. In multiplayer: arrivare primo, e soprattutto
\textbf{mostrare} i propri ottenimenti. Nella relazione tra achiever e socializer,
di cui parla lo stesso Bartle, conta tutto ciò che l'achiever ha ottenuto e può
esibire: il banner, il mantello esclusivo del dicembre 2012, la skin unica del
carnevale di Rio 2002, ``questa roba una tantum che non può più essere ottenuta''.
Per l'achiever è importante farla vedere; la slide aggiunge che si rivolge ai
socializer per riceverne lodi.

\paragraph{Explorer.} In single player: tutto ciò che è \textbf{osservabile,
tracciabile e narrabile, coerentemente}; il \emph{backstory}, la \emph{lore}. Qui c'è
un altro punto di contatto, ``molto pericoloso'': come quello tra killer e achiever è
il diventare numero uno, quello tra explorer e socializer è \textbf{la storia, la
narrativa}. Si distingue così: se si parla di qualcosa di altamente regolamentato,
della lore, è explorer. Molti giocatori di \emph{Dark Souls} uniscono i puntini della
lore come davanti alla classica lavagna investigativa (``questo è capitato prima,
questo è capitato dopo''): il presupposto è che il mondo sia \textbf{coerente}, e
l'obiettivo è muoversi entro le sue regole. Il godere della trama, del
\textbf{dramma}, invece, ricade nel socializer: il dramma sociale, la storia per il
gusto di vedere qualcosa di narrato. In multiplayer l'explorer cerca nuovi contenuti
e \textbf{guida gli altri} verso i segreti, condivide le informazioni, usa gli altri
giocatori come perno (non è detto: potrebbe anche essere a cavallo tra killer ed
explorer).

\paragraph{Socializer.} Il multiplayer è il suo ambiente naturale: come per il killer
il multiplayer offre la preda definitiva, un giocatore imprevedibile, per il
socializer offre \textbf{l'amico definitivo}, una persona che risponde socialmente
ai suoi stimoli e viceversa, con un potenziale quasi illimitato di nuove relazioni.
In single player cerca l'esperienza di comunità (forum, wiki) o la
\textbf{simulazione di relazioni}: non necessariamente il dramma in sé, ma
spiccatamente il \textbf{dramma sociale}. \emph{The Last of Us} o il nuovo \emph{God
of War}, in cui si gioca nei panni di una persona in relazione con un altro NPC,
``figlio adottivo o meno'', sono spiccatamente da socializer. E anche \emph{The Sims}
o \emph{Tomodachi Life}, in cui si costruiscono relazioni sociali: ``più persone,
meglio è''.

\subsection{Una domanda: a che cosa serve?}

Uno studente chiede, per avere chiaro dove si sta andando: la tassonomia definisce
un giocatore rispetto a un determinato gioco, o il modo in cui un giocatore
approccia i giochi in generale? ``Possono essere entrambe le cose''; l'obiettivo
finale si può usare in diverse maniere, e lo si vedrà alla fine della lezione. Ma già
ora, con i dati presi in pre-produzione, si può individuare \textbf{dove è
posizionato il proprio target audience} (non necessariamente tutti killer: una
percentuale di killer, una di explorer\dots), e anche dove sono posizionati i
\textbf{competitor}: quali meccaniche hanno, come vengono usate e a quale player type
sono dirette. Questo dà già una direzione. Un altro modo d'uso è la soluzione dei
problemi: davanti a un problema che richiede una soluzione creativa, ci si chiede
\textbf{a quale di questi tipi piacerebbe}. Se il target è l'achiever, che cerca una
lista di cose da fare, probabilmente un \textbf{quest log esplicito} gli sarà molto
utile.

Il docente chiede un minuto per andare a riempire la borraccia, ``giusto per non
morire''.\footnote{La breve pausa è riempita, nella trascrizione, dalle solite frasi
in loop; al ritorno qualcosa si blocca nella condivisione (``siamo fregati\dots
completamente bloccato''), poi la lezione riprende.}

% ---------------------------------------------------------------------
\section{I temperamenti di Keirsey}
\label{gd18:sec:keirsey}

\slidefig{img/GD/gd18-keirsey}{I quattro temperamenti di Keirsey.}{slide \emph{DM06 -- Persona \& Types}}{gd18:fig:keirsey}

\subsection{Gli assi di Bartle, cambiati}

Un altro modo di vedere la questione, e normalmente sarebbero modelli distinti:
\textbf{David Keirsey} parte dagli assi di Bartle e li cambia, cercando di arrivare
allo stesso risultato (figura~\ref{gd18:fig:keirsey}).\footnote{David Keirsey,
\emph{Please Understand Me II: Temperament, Character, Intelligence}, 1998. Nella
trascrizione il nome è reso come ``Chirsey'', ``Kirsey'' o ``Kirsten''.}
\begin{itemize}
  \item invece di \emph{acting on} c'è \textbf{esterno} (\emph{externals});
  \item invece di \emph{interacting with} c'è \textbf{interno} (\emph{internals});
  \item invece di \emph{players} c'è \textbf{cambiamento} (\emph{change});
  \item invece di \emph{world} c'è \textbf{struttura} (\emph{structure}).
\end{itemize}
Questo rende i due framework \textbf{complementari}, e soprattutto ha portato
all'\textbf{espansione}: a cercare quali altre visioni e modelli si potessero
integrare, appoggiandosi alla tassonomia di Bartle per avere uno spettro più ampio
di definizione dei target audience.

\subsection{I quattro temperamenti}

\paragraph{Artisan (SP, \emph{sensing + perceiving}).} Corrisponde al killer.
Cambiamento \textbf{esteriore}: vuole apportare un cambiamento al di fuori di sé.
Cerca le \textbf{sensazioni}: l'adrenalina, la ricerca ``spasmodica'' del picco
emotivo, le grandi emozioni, l'\emph{excitement}. L'\emph{action} di Quantic Foundry
è del killer, ``non unicamente, ma principalmente''. È realistico, \textbf{tattico}
(a breve termine), manipolativo, focalizzato sull'azione e sulle sensazioni.

\paragraph{Guardian (SJ, \emph{sensing + judging}).} Corrisponde all'achiever.
Cerca una \textbf{struttura esteriore}, cioè le leggi e i limiti del gioco, e gioca
entro i limiti imposti. Pratico, logistico, organizzato, attento al dettaglio,
possessivo, concentrato sui processi; la sua ricerca principale è la
\textbf{sicurezza}, ottenuta raggiungendo tutti gli obiettivi, ragionando dentro i
limiti del gioco, accumulando risorse. L'obiettivo finale dell'achiever è
\textbf{avere}; quello del killer, \textbf{fare}.\footnote{Il docente, a voce,
contrappone il guardian ``strategico'' all'artisan ``tattico''; nelle slide e nella
tabella del modello prismatico (tabella~\ref{gd18:tab:prismatico}) il guardian è
\emph{logistical} e lo \emph{strategic} è il rational.}

\paragraph{Rational (NT, \emph{intuition + thinking}).} Corrisponde all'explorer.
\textbf{Struttura interna}: vuole avere chiara, mentalmente, la struttura del gioco.
Innovativo, logico, scientifico, orientato al futuro, focalizzato sui risultati;
cerca la \textbf{conoscenza} (\emph{knowledge}). Non a caso, una cosa spiccatamente
apprezzata dagli explorer sono i \textbf{puzzle}, perché un puzzle è generalmente un
\textbf{problema di percezione}.

\paragraph{Idealist (NF, \emph{intuition + feeling}).} Corrisponde al socializer. Il
suo obiettivo è \textbf{modificare sé stesso}: attraverso la relazione con gli
altri scopre qualcosa di sé, cresce, evolve. Immaginativo, \textbf{diplomatico} (se
l'artisan cerca sempre il conflitto, l'idealist cerca la diplomazia, la pace),
emotivo, orientato alle relazioni, concentrato sulle persone; cerca
l'\textbf{identità}, in sé stesso e negli altri: scoprire gli altri è il metodo, e il
metro, per capire sé stessi.

I temperamenti di Keirsey si sovrappongono uno a uno con i tipi di Bartle, ``anche
perché Keirsey lì arriva''.

\subsection{Le sedici personalità}

Prima di andare avanti, il docente allinea un altro elemento che fa già parte del
framework: le \textbf{sedici personalità}.\footnote{Le sedici combinazioni del
\emph{Myers-Briggs Type Indicator} (MBTI), su cui si basano i temperamenti di
Keirsey: ogni temperamento ne raccoglie quattro.} Non sono nate per questo, ma
``fortunatamente'' sono divisibili per quattro, e si possono inserire nella
tassonomia di Bartle. Gli studenti dovrebbero averle fatte durante il colloquio; il
docente è \textbf{Advocate}.

\slidefig{img/GD/gd18-bartle-keirsey}{Bartle e Keirsey sovrapposti: gli assi uniti e le quattro combinazioni.}{slide \emph{DM06 -- Persona \& Types}}{gd18:fig:bartle-keirsey}

\subsection{Bartle-Keirsey}

Sovrapponendo i due modelli (figura~\ref{gd18:fig:bartle-keirsey}), con gli assi di
Bartle sotto e quelli di Keirsey sopra, emergono i punti in comune:
\begin{description}
  \item[Artisan/Killer.] Manipolatori, che vogliono solo dimostrare la propria
        \textbf{superiorità} sugli altri esseri umani, e in ultimo in generale:
        arrivare primi, avere il tempo migliore, uccidere tutti, mettere i bastoni tra
        le ruote, essere l'ultimo vivo.
  \item[Guardian/Achiever.] Il mondo è un posto molto \textbf{insicuro}, quindi la
        sicurezza passa per l'\textbf{accumulo materiale}, potenzialmente anche
        competendo con gli altri per risorse percepite come scarse. E una risorsa può
        essere la prima posizione, che è una sola; in qualsiasi multiplayer la prima
        posizione è legata a ottenimenti materiali (il banner della stagione, i soldi,
        la valuta di gioco). Ma anche gli achievement: completare un gioco, per
        alcuni, è il 100\% degli achievement.
  \item[Rational/Explorer.] Cercano l'\textbf{ordine dietro i dati grezzi}: capire
        che cosa c'è dietro il velo del gioco, conoscere le meccaniche, i sistemi, come
        si relazionano tra loro, il sistema come un tutto.
  \item[Idealist/Socializer.] Vogliono giocare con l'altra persona, stare insieme
        agli altri, capirli e, in ultimo, creare \textbf{relazioni soddisfacenti} e
        durature.
\end{description}
Il docente chiede se è tutto chiaro; manca ancora un pezzo. Finora si è parlato
degli \textbf{estremi}, mai delle vie di mezzo: chi è a metà tra explorer e
socializer, o tra socializer e achiever?

% ---------------------------------------------------------------------
\section{Il modello demografico di Bateman}
\label{gd18:sec:bateman}

\slidefig{img/GD/gd18-bateman}{Il modello demografico DGD1 di Chris Bateman, con i tipi intermedi.}{slide \emph{DM06 -- Persona \& Types}}{gd18:fig:bateman}

\subsection{I tipi intermedi}

I player type ``a metà'' sono quelli di \textbf{Chris Bateman}, autore di
\emph{21st Century Game Design}, il libro che il docente aveva consigliato
(capitolo~\ref{cap:introduzione-game-design}). Il suo \textbf{Demographic Game
Design Model} (DGD1, figura~\ref{gd18:fig:bateman}) colloca sul quadrato di
Bartle-Keirsey i tipi intermedi:
\begin{description}
  \item[Conqueror] (strategico-logistico, tra rational/explorer e guardian/achiever):
        associato alla \textbf{sfida} e alla ricompensa emotiva del \emph{fiero}; è
        molto paziente con la frustrazione, perché sa che, se insiste, può vincere.
  \item[Manager] (strategico-tattico): associato alla \textbf{mastery} e ai
        \textbf{sistemi}. Per chi preferisce questo stile la vittoria sembra il segno
        di aver acquisito le abilità necessarie, non un fine in sé; è bravo a gestire
        più fattori in parallelo.
  \item[Wanderer] (diplomatico-tattico, tra idealist/socializer e
        artisan/killer):
        associato all'\textbf{esperienza} e all'\textbf{identità}. La sfida non è
        particolarmente desiderata: l'obiettivo è godersi esperienze uniche e
        interessanti.
  \item[Participant] (diplomatico-logistico): associato alle \textbf{emozioni} e al
        \textbf{coinvolgimento}. Sembra più felice quando gioca con altre persone, ma
        apprezza anche giochi radicati nell'emozione.
\end{description}
Lo schema aggiunge due tipi lungo i lati: \textbf{Hardcore}
(diplomatico-strategico, sul lato \emph{internals}) e \textbf{Casual}
(tattico-logistico, sul lato \emph{externals}), e accanto a ogni tipo le quattro
sigle delle sedici personalità che vi ricadono.

\subsection{Un disclaimer: il casual di oggi}

Il grande vantaggio di Bateman, e il motivo per cui il docente lo mostra, è
\textbf{aprire la mente ai tipi intermedi}, che è il dubbio che forse stava già
emergendo in qualcuno. Ma alcuni tipi sono ``vecchiotti, con tutto rispetto''.
L'esempio è proprio il \textbf{casual}, a cavallo tra killer e achiever,
tattico-logistico: non è il casual gamer del 2026; deriva da un mercato e da un tipo
di giocatore che non c'è più (o meglio, c'è, ma è un altro tipo di giocatore).

Chi sarebbe il casual gamer di oggi? Il docente lo chiede alla classe. ``Mia zia che
gioca a \emph{Clash of Clans}'', risponde uno studente. ``Certo, ma tua zia che gioca
a \emph{Clash of Clans} difficilmente la vedo come killer-achiever.'' Bisogna capire
dove si colloca: forse socializer-achiever, forse explorer-achiever, ``sicuramente
non killer''. E chi gioca a \emph{FIFA}? È un altro discorso: i giocatori di
\emph{FIFA} starebbero forse lì, a seconda dell'obiettivo. Il casual di Bateman è
tattico-logistico perché la stragrande maggioranza del mercato, a quel tempo, era
fatta di giocatori di \textbf{FPS}, che sono esattamente killer-achiever:
\emph{Gears of War}, \emph{Call of Duty}, \emph{Halo}, ``siamo esattamente qua''.

Anche gli altri vanno letti con cautela. Il \textbf{Wanderer}, o il \emph{lone wolf}
secondo un altro framework, è una persona in solitario: si pensi a Gandalf, ``a metà
tra grigio e bianco'', che fa le sue camminate da solo e ogni tanto fa parte della
Compagnia, ogni tanto no. E anche l'\textbf{hardcore} non è più ciò che era: forse
oggi starebbe altrove. L'obiettivo principale, ripete il docente, è avere chiaro in
testa che si tratta di uno \textbf{spettro}: ci sono persone con più o meno
predilezione per un archetipo, ma a cavallo tra due, o tra tre.

\paragraph{Approfondimento: hardcore e casual nel DGD1.} Nel libro di Bateman e
Boon,\footnote{Chris Bateman, Richard Boon, \emph{21st Century Game Design}, 2005,
capitolo 4.} i quattro tipi del DGD1 nascono incrociando i dati sulle preferenze di
gioco con le sedici personalità: Conqueror (TJ: ISTJ, INTJ, ESTJ, ENTJ), Manager (TP),
Wanderer (FP) e Participant (FJ). Gli autori avvertono che vanno letti come un
\emph{continuum} e non come caselle rigide, e che non si escludono a vicenda. La
distinzione tra hardcore e casual è definita per comportamenti: l'hardcore compra e
gioca molti giochi, ama sessioni lunghe, sfida, progressione e mastery, tollera
controlli complessi e vive il gioco come uno stile di vita; il casual compra meno
giochi (quelli popolari o consigliati dagli amici hardcore), gioca a sessioni brevi,
preferisce divertirsi o immergersi in un'atmosfera, vuole controlli semplici e vede
il gioco come un passatempo, come la TV o i film. È la definizione che, vent'anni
dopo, il docente invita ad aggiornare.

% ---------------------------------------------------------------------
\section{Il modello prismatico}
\label{gd18:sec:prismatico}

\subsection{Allargare Bartle, con giudizio}

Bartle, Keirsey e Bateman (soprattutto gli ultimi due) hanno costruito sulle
fondamenta di Bartle. Ma il modello di Bartle ha un elemento fondante
\textbf{indiscutibile}: gli archetipi sono \textbf{quattro}. È il motivo per cui le
sedici personalità si sposano così bene, perché sono divisibili per quattro; ma non
vale per tutto. L'intento è nobile (allargare il framework per avere più punti di
vista sullo stesso player type), ma non sempre è possibile, e a volte è
\textbf{forzato}.

Qui ``rientra il come volete vederla voi''. Uno strumento in più è probabilmente
molto utile a un designer, a maggior ragione se si sposa con ciò che vuole fare e
con ciò che sente; ma deve essere \textbf{attinente} a ciò che fa e a come lo usa. Da
qui in avanti si vedono altri modi di integrare la tassonomia di Bartle, che portano
al modello unificato, o prismatico, ``potete chiamarlo come volete'': unire più cose
per avere una visione più ampia. Ciò \textbf{non} significa sposarle tutte. Molti di
questi modelli vengono dalla psicologia, addirittura dall'economia o dalla gestione
aziendale (Gallup, Covey\dots); alcune definizioni, per un singolo archetipo, sono
molto vantaggiose, altre meno. Dipende da quante cose si vogliono inserire e da
quanta ampiezza si vuole dare al framework. Un punto, però, non è negoziabile:
\begin{quote}
Voi partite da Bartle e dovete saperlo a memoria. [\dots] Quanto espanderlo, cosa
utilizzare, rimane a voi. Tuttavia, prima, come sempre, cercate di masterare le
basi.
\end{quote}

\slidefig{img/GD/gd18-caillois}{Le quattro categorie di gioco di Caillois: agon, alea, mimicry, ilinx.}{slide \emph{DM06 -- Persona \& Types}}{gd18:fig:caillois}

\subsection{Caillois: agon, alea, mimicry, ilinx}

Nel modello unificato si è andati a prendere anche il ``buon vecchio''
\textbf{Roger Caillois}, francese del secolo scorso (il docente esita sul nome, poi
lo trova): a dimostrazione che si possono prendere teorie e scritti anche
\textbf{precedenti} a Bartle e incastrarli nel framework. Caillois è l'autore de
\emph{I giochi e gli uomini},\footnote{Roger Caillois, \emph{Les jeux et les
hommes}, 1958.} che prima distingue il gioco regolato e il gioco sregolato
(\emph{ludus} e \emph{paidia}, ``non facciamo il discorso generale''), e poi divide
la ricerca e la motivazione nel gioco in quattro categorie
(figura~\ref{gd18:fig:caillois}). Al docente Caillois piace particolarmente, e
``sposa molto la causa'' di integrare queste quattro categorie; ma potrebbe non
valere per tutti, anche perché Caillois tratta gli \textbf{estremi}.

\begin{description}
  \item[Agon.] Lo sforzo fisico, tutto ciò che è \textbf{agonistico}: pugilato,
        calcio, basket. Il senso agonistico del raggiungere il primato con tanta
        fatica, la mastery dei sistemi, la strategia. L'esempio sono gli
        \textbf{scacchi}: un gioco che si dice simmetrico (anche se non esattamente,
        perché uno muove prima dell'altro) perché le risorse sono uguali per
        entrambi. Non si possono potenziare o personalizzare i pezzi (l'unica
        eccezione è la promozione); non si può decidere prima della partita di
        giocare con tre cavalli mentre l'altro gioca con quattro torri; non si
        draftano i pezzi né si pesca una carta a caso. È un gioco spiccatamente
        agonistico.
  \item[Alea.] L'\textbf{aleatorietà}: tutto ciò che è imprevedibile, la fortuna, il
        pescare, il lancio dei dadi, la lotteria. Caillois ragiona anche sugli
        estremi e su come la società integra questi elementi: l'alea portata
        all'estremo, la sua disfunzione, è la \textbf{ludopatia}.
  \item[Mimicry (o mimesis).] \textbf{Mettere una maschera}: il
        \emph{make-believe}, l'impersonazione di qualcos'altro, regolata. GDR,
        teatro, giochi di finzione. La disfunzione è la dissociazione della
        personalità: non necessariamente personalità multiple, ma il mettere
        costantemente una maschera, avere maschere sociali molto diverse in contesti
        diversi.
  \item[Ilinx.] La \textbf{vertigine}: voler cercare emozioni forti, il rischio per
        la vita, gli sport estremi, il funambolo. L'estremo è la dipendenza da
        droghe, cercare l'alterazione del proprio stato, ubriacarsi: la
        disinibizione da qualsiasi limite (``vabbè, non entriamo nei dettagli'').
\end{description}

\paragraph{Il D\&D del docente.} Per spiegare la mimicry il docente porta il proprio
caso. Lui, a differenza di altri colleghi, è spiccatamente \textbf{meccanico}: gli
piacciono molto le meccaniche dei giochi, mentre la parte narrativa, il
make-believe, ``non dico che non sia importante, ma è meno importante per me''.
Gioca a \emph{Dungeons \& Dragons} ragionando su ``colpisce, non colpisce, ho
abbastanza vita, non ho abbastanza vita'', e i suoi personaggi sono sempre ``il nano
incazzato che deve menare'', in qualsiasi ambientazione, cyberpunk o fantasy. Ma non
è così che si gioca davvero: a D\&D si impersona un personaggio diverso da ciò che
si è fuori dal tavolo, ``l'elfo Gilendal con un determinato passato alle spalle'' che
agisce in una determinata maniera. Questo è la mimicry.

\subsection{Il quiz: a chi va ciascuna categoria?}

Il docente chiede alla classe di abbinare le quattro categorie ai player type.
\textbf{Agon} e \textbf{ilinx} sono ``facili'':
\begin{itemize}
  \item l'\textbf{agon} è dell'\textbf{achiever}, per quel discorso di sposare le
        regole: il primato raggiunto con fatica, la mastery dei sistemi, la strategia
        agonistica. Uno studente obietta con il \textbf{calcio}: è uno sport di
        squadra, e come altri sport non è indirizzato a un unico player type;
  \item l'\textbf{ilinx} è del \textbf{killer}. Uno studente propone il
        \textbf{paracadutismo}: se si riesce a soddisfare il proprio obiettivo, non è
        da achiever? La differenza, risponde il docente, è che l'achiever
        \textbf{non si metterebbe mai in pericolo di vita}.
\end{itemize}
Ci sono, certo, accezioni comuni e intermedi; per ora si cerca di tenere distinte le
quattro categorie. Tattica e strategia, per esempio, sono comuni a killer e
achiever. E c'è una cosa che sta proprio a metà tra i due: la
\textbf{challenge}. Alcune emozioni (\emph{action}, \emph{excitement},
\emph{curiosity}, \emph{discovery}\dots, tutto ciò che si è visto con Quantic
Foundry) sono più apprezzate da alcuni player type che da altri, e a volte da più
d'uno. La domanda da porsi, come detto a Lorenzo,\footnote{Probabilmente in
riferimento alla domanda sui cheat (sezione~\ref{gd18:sec:bartle}).} è
\textbf{perché} il player type vuole qualcosa, qual è il suo obiettivo: sulla base
di ciò che desidera lo si inquadra meglio. Imporsi su una sfida, diventare il più
forte, ``distruggere il boss di legnate'', è da killer; superare una grossa sfida,
battere il boss tramite mastery o perché si è grindato, è da achiever. La sfida in
sé (boss, sfide a tempo, challenge con restrizioni) sta a cavallo tra i due: va
capito non quante meccaniche ha un sistema, ma \textbf{come il giocatore le prende in
mano}.

\paragraph{Mimicry e alea.} Le altre due sono ``più toste''. Uno studente propone la
mimicry per il socializer, perché ricerca l'attenzione delle persone impersonando
qualcuno che non è; per l'alea si azzarda l'explorer, per la curiosità verso ciò che
è imprevedibile (ma un altro obietta che non poter prevedere il risultato metterebbe
in difficoltà l'explorer). Il docente dà la risposta: vanno \textbf{invertite}.
\begin{center}
\textbf{agon} $\rightarrow$ achiever \qquad \textbf{alea} $\rightarrow$ socializer
\qquad \textbf{mimicry} $\rightarrow$ explorer \qquad \textbf{ilinx} $\rightarrow$
killer
\end{center}

\paragraph{Perché la mimicry è dell'explorer.} ``Ci siamo quasi'', dice il docente a
chi parla di cambiare maschere. Il make-believe \textbf{ha delle regole}: se si
impersona un alto elfo, certe cose non le si farà, o non sarà corretto farle; non ci
si vorrà sporcare, non si vorranno armi da mischia, non si diranno certe cose e se
ne faranno altre. Avere diverse maschere, ma con regole specifiche, ricade in quel
\textbf{unire i puntini della lore}, capire coerentemente la storia: la parte
prevedibile, regolata, è dell'explorer. È per questo che a lui Caillois piace tanto:
offre un altro punto di vista sui player type.

\paragraph{Perché l'alea è del socializer.} Quanti hanno mai giocato a
\emph{Exploding Kittens}? E perché \emph{Mario Kart} funziona così bene in compagnia
(``ancora meno gli ultimi \emph{Mario Kart}'')? L'alea non solo crea molta
\textbf{speranza di ribaltare la situazione}, ma \textbf{ci rende tutti pari}:
\begin{quote}
Davanti alla sorte noi siamo uguali. Quindi è indifferente che tu abbia una mastery
enorme dei sistemi, o che abbia cento milioni di livelli, o che sia il più forte
giocatore di \emph{Call of Duty} al mondo: se è modalità random, arma random, e io mi
becco l'RPG e tu il coltello, è chiaro che tu hai uno svantaggio enorme.
\end{quote}
Funziona bene anche nei giochi di carte, sia nell'apertura dei pacchetti sia nello
sfidare un avversario: se si è al livello 1 e dall'altra parte c'è il miglior
giocatore del mondo al livello 50, l'unica speranza è sopraffarlo con la fortuna. Lo
stesso vale per il poker, o per un semplice lancio di dadi. Un gioco completamente
aleatorio, come il gioco dell'oca o \emph{Scale e serpenti}, è chiaramente pensato
per i bambini e ha un pubblico molto giovane.

Non si può però basare qualsiasi gioco di carte sull'alea: molti sono spiccatamente
da achiever (\emph{Hearthstone}, \emph{Yu-Gi-Oh!}, il gioco di carte Pokémon). In
quei giochi si inseriscono altre meccaniche per \textbf{ridurre la magnitudine
dell'alea}: se si pesca male a \emph{Yu-Gi-Oh!} si perde quella partita, e per questo
di solito si gioca al meglio delle tre; a \emph{Backgammon} i lanci possono andare
male, ma il giocatore davvero abile riesce a sopraffare quello meno abile anche con
un grande scarto di fortuna, perché lì entra la mastery del sistema.

\paragraph{Il poker e il novizio.} Il poker, accenna il docente, è un altro gioco
interessante: la strategia migliore per un novizio a un tavolo di esperti è
\textbf{non guardare le proprie carte}. Così non dà nessuna informazione agli altri
su ciò che ha visto o non ha visto, e sbilancia il tavolo verso quella parità
dell'alea, ``perché nemmeno tu sai cos'hai'': giocando a caso, probabilmente vince
molte più mani di quante ne vincerebbe giocando.

\slidefig{img/GD/gd18-lazzaro}{Le quattro chiavi del divertimento di Nicole Lazzaro.}{slide \emph{DM06 -- Persona \& Types}}{gd18:fig:lazzaro}

\subsection{Lazzaro: le quattro chiavi del divertimento}

Poi c'è \textbf{Nicole Lazzaro}, ``l'unico modo di dire \emph{fun} senza essere
bacchettati'' (figura~\ref{gd18:fig:lazzaro}). Le sue \textbf{quattro chiavi del
divertimento} (già incontrate nel capitolo~\ref{cap:ricerca-idee}) sono molto utili,
anche per un motivo di metodo: molti dati, come quelli di Bartle, \textbf{non sono
pubblici}; Lazzaro invece ha fatto un survey, ha osservato i giocatori, li ha
registrati e ha cercato di categorizzarne le risposte emotive, anche facciali. Il
risultato descrive alcune risposte molto legate ai tipi di giocatore. E c'è un
dettaglio importante: se il \textbf{fiero} non può essere allineato a nessuno dei
player type visti finora (almeno non a uno solo), l'\textbf{hard fun} ci va molto
vicino: è ``la cosa più vicina al fiero che potete dire senza dire fiero''.

\begin{description}
  \item[Serious fun.] Un \textbf{picco di adrenalina} seguito da un \textbf{rilascio
        emotivo}: \emph{Alien: Isolation}, tutti i giochi horror, in cui una sezione
        estremamente adrenalinica è seguita da una di rilascio. È chiaramente del
        \textbf{killer}: l'adrenalina, il picco emotivo, la vertigine.\footnote{La
        slide riporta la definizione più recente di Lazzaro, secondo cui il
        \emph{serious fun} è il gioco con uno scopo, che cambia il modo in cui i
        giocatori pensano, sentono e si comportano, e che rende vivi compiti
        altrimenti noiosi (``play as therapy''). La lettura del docente, in termini di
        eccitazione e rilascio, è più vicina alla versione originale della chiave,
        chiamata \emph{altered states} (si veda anche il
        capitolo~\ref{cap:ricerca-idee}).}
  \item[Hard fun.] Il \textbf{superamento di una sfida} tramite la mastery dei
        sistemi: una sfida molto ardua, si è perso tante volte, e finalmente arriva
        il sospiro di sollievo, la gioia della vittoria. C'è un obiettivo, ci sono dei
        vincoli e c'è una strategia per superarli. È dell'\textbf{achiever}.
  \item[Easy fun.] La parola che lo descrive meglio è \textbf{curiosità}: autonomia,
        indagine, esplorazione, roleplay, fantasia. È ciò che si prova al primo
        approccio con un puzzle: la sperimentazione, l'approccio giocoso a qualcosa
        che si vede per la prima volta, il curiosarci, esplorarlo. È
        dell'\textbf{explorer}.
  \item[People fun.] Comunità, divertimento, la sensazione di tranquillità e di
        piacere di stare in un gruppo di amici con cui si può stare tranquilli: i
        party game, la scusa per stare con gli altri o parlarci, il \emph{social
        bonding}. È del \textbf{socializer}.
\end{description}

Di modelli ce ne sono molti altri; il docente ha mostrato i più importanti, e
precisa la gerarchia: ``per quanto mi piaccia Caillois'', è \textbf{meno importante
di Lazzaro}. La priorità va a Lazzaro.

\slidefig{img/GD/gd18-prismatico}{Il modello prismatico (o unificato): i modelli allineati sui quattro tipi di Bartle.}{slide \emph{DM06 -- Persona \& Types}}{gd18:fig:prismatico}

\subsection{La tabella del modello prismatico}

Il modello unificato in partenza, poi chiamato prismatico da Diego, si basa su
questa \textbf{ripartizione a quattro} (figura~\ref{gd18:fig:prismatico} e
tabella~\ref{gd18:tab:prismatico}) e lo si può ampliare quanto si vuole.\footnote{La
tabella riprende quella dell'articolo di Bart Stewart, \emph{Personality and Play
Styles: A Unified Model}, pubblicato su Gamasutra nel 2011, che allinea Keirsey,
Bartle, Caillois, Lazzaro, GNS+, MDA+, Handy, Gallup e Covey sui quattro stili di
gioco.}

\begin{table}[ht]
\centering
\scriptsize
\setlength{\tabcolsep}{3pt}
\resizebox{\textwidth}{!}{%
\begin{tabular}{@{}llllllllll@{}}
\toprule
\textbf{Keirsey} & \textbf{Bartle} & \textbf{Caillois} & \textbf{Lazzaro} &
\textbf{Handy} & \textbf{Gallup} & \textbf{Covey} & \textbf{Motivazione} &
\textbf{Problem solving} & \textbf{Obiettivo} \\
\midrule
Artisan (tattico) & Killer & ilinx & serious fun & Power & Impacting & Power &
Power & Performance & Do (fare) \\
Guardian (logistico) & Achiever & agon & hard fun & Role & Striving & Security &
Security & Persistence & Have (avere) \\
Rational (strategico) & Explorer & mimesis & easy fun & Task & Thinking & Wisdom &
Knowledge & Perception & Know (conoscere) \\
Idealist (diplomatico) & Socializer & alea & people fun & People & Relating &
Guidance & Identity & Persuasion & Become (diventare) \\
\bottomrule
\end{tabular}}
\caption{Il modello prismatico, senza le colonne GNS+ e MDA+ (che il docente chiede
di non usare). Motivazioni per esteso: \emph{power} (sensazione manipolativa),
\emph{security} (accumulo competitivo), \emph{knowledge} (scoperta logica delle
regole), \emph{identity} (relazioni emotive).}
\label{gd18:tab:prismatico}
\end{table}

\subsection{Due colonne da non usare: GNS+ e MDA+}

Non bisogna sposare tutte le cause, e il docente dice la sua, anche se ``non è che
si dica spesso'': alcune colonne sono, a suo avviso, \textbf{sbagliate} o comunque
non applicabili, perché nascono da un processo di \textbf{forzatura}.
\begin{itemize}
  \item \textbf{GNS+}. È un framework che forse gli studenti vedranno con Daniele (o
        forse no, perché le sue lezioni sono poche), legato in qualche modo al
        \textbf{cerchio magico di Huizinga} (``non mi ricordo se è venuto prima
        l'uovo o la gallina'') e alla narrazione nel gioco, temi che verranno
        affrontati nelle lezioni di narrativa. Si fonda su \textbf{gamist},
        \textbf{simulationist} e \textbf{narrativist}: questo è il pezzo importante.
        L'\textbf{experientialist} ``non esiste'': è stato aggiunto per far tornare
        il conto a quattro. Al docente il GNS di base non fa impazzire; per lui può
        allinearsi, ma ``scegliete quello che vi pare''.
  \item \textbf{MDA+}. Qui il docente è netto: ``per il vostro benessere psicofisico
        all'interno del corso, lo togliete da qui''. Per come verrà trattato nel corso
        e per come lo approcceranno anche gli altri docenti, il punto è sempre lo
        stesso: le \textbf{kinetics} non esistono nel framework MDA di base, e anche
        se lo si volesse espandere (il design è un processo iterativo, si va sempre
        verso il meglio), \textbf{non è l'MDA che va integrato in questo framework,
        è il contrario}. È l'MDA che si può appoggiare al modello prismatico. Lo si
        capirà meglio quando si farà l'MDA l'anno prossimo; per ora ``tenetelo nella
        saccoccia''. Ci si baserà su ciò che dice questo framework (e sulle sue
        espansioni) per ottenere dall'MDA \textbf{dati quantitativi}.
\end{itemize}

\subsection{Le altre colonne}

Ci sono poi modelli che non sono legati al design di per sé.
\begin{description}
  \item[Handy.] Charles Handy parla di \textbf{cultura aziendale}: \emph{power
        culture} (risolvere i problemi appena arrivano, distruggendo le gerarchie in
        funzione di una soluzione veloce: un certo dinamismo aziendale); \emph{role
        culture} (la gerarchia netta e definita, eventualmente piramidale, e restare
        nei suoi limiti); \emph{task culture} (avere tante cose da fare e portarle
        avanti contemporaneamente); \emph{people culture} (seguire leader
        carismatici, un po' come il lone wolf visto prima, a capo di un gruppo di
        socializer).
  \item[Gallup.] \textbf{Impacting} per il killer (che è proprio il suo obiettivo:
        impattare gli altri giocatori); \textbf{striving} per l'achiever (mirare a
        qualcosa, sforzarsi di ottenere una spada, una moneta, un achievement);
        \textbf{thinking} per l'explorer; \textbf{relating} per il socializer.
  \item[Covey.] Gli obiettivi: \textbf{power} (ottenere potere), \textbf{security},
        \textbf{wisdom}, e per il socializer \textbf{guidance}, cioè guidare gli altri
        e sé stesso verso la conoscenza personale, ``mia e dell'altro''.\footnote{A
        voce il docente assegna \emph{power} all'achiever e \emph{security} a un
        secondo tipo, correggendosi più volte; la tabella delle slide è quella
        riportata nella tabella~\ref{gd18:tab:prismatico}.}
  \item[Motivazione.] Power, security, knowledge, identity: cose già dette in alcune
        definizioni.
  \item[Problem solving.] Come risolvono i problemi i singoli tipi: il killer
        tramite la \textbf{performance}; l'achiever tramite la \textbf{persistenza}
        (il grinding, per lui, non è solo utile, è anche piacevole); l'explorer
        tramite la \textbf{percezione}, cambiando punto di vista (ed è così che di
        solito si risolvono i puzzle); il socializer tramite la
        \textbf{persuasione}, con la comunicazione.
  \item[Obiettivo finale.] Il killer vuole \textbf{fare}, l'achiever \textbf{avere},
        l'explorer \textbf{conoscere}, il socializer \textbf{diventare}.
\end{description}

% ---------------------------------------------------------------------
\section{Il ``non task'': rifarsi il modello unificato}
\label{gd18:sec:rifare}

``Io non vi do un task per la prossima lezione, ve lo dico già, così potete tirare un
respiro di sollievo.'' Però c'è qualcosa da fare, che ``forse potreste vedere come un
task'': \textbf{rifarsi questo modello unificato}. Da qui in avanti il docente si
aspetta che lo si usi non solo nel processo di design, ma integrato in qualsiasi task
venga dato: deve far parte della quotidianità del designer.

Il consiglio:
\begin{itemize}
  \item rifarselo come tabella, o non come tabella; il docente manderà un altro
        documento che integra un po' di cose e può servire da base;
  \item fare \textbf{ricerca sui singoli modelli}: non si è trattato tutto di Handy,
        di Gallup, di Covey e così via;
  \item lasciar perdere GNS+ (``sconsigliato, fate come vi pare'') e soprattutto MDA+
        (``proibisco l'MDA''). Anche della \emph{dynamic difficulty adjustment} si
        parlerà più avanti;
  \item ricordare soprattutto le componenti di base di \textbf{Bartle, Keirsey e
        Lazzaro}: ``li dovete sapere a memoria'', e questo è un modo per aiutarsi a
        farlo.
\end{itemize}
Il documento che ne verrà fuori (``un bigino, come volete chiamarlo, un documento di
ritorno, un benchmark'') servirà anche a capire che cosa piace o non piace: piace
Lazzaro? Piace come divide Caillois? ``Bartle no, guarda che vi deve piacere per
forza, perché sono i termini.'' L'unico modo è conoscerli.

\paragraph{Le case di Hogwarts.} Federica interviene: nella sua testa era partito
come uno scherzo, ma i player type si potrebbero dividere anche nelle quattro
\textbf{case di Hogwarts}. ``Tecnicamente funziona, certo. Sorprendentemente sì.''
Nel documento che il docente manderà, una vecchia versione di Diego da cui partire,
ci sono per esempio anche le \textbf{classi}: il mago, il ladro (``chiaramente il
ladro è il killer''). Il docente non le usa, così come non usa le
\textbf{personificazioni} (le mani, gli occhi, lo spirito\dots): sono difficili da
spiegare e non fanno parte del suo processo di design. Ma potrebbero fare parte di
quello degli studenti, aiutarli a ragionare meglio o dare nuovi spunti su come
vedere un player type. L'obiettivo è integrare in un framework tutto ciò che si sa e
che si vuole usare: ``perché da questo non si sfugge, mi spiace''.

% ---------------------------------------------------------------------
\section{L'analisi per generi}
\label{gd18:sec:generi}

Il passo successivo, forse oggetto della lezione seguente: rispetto ai propri giochi
preferiti, \textbf{identificare i player type principali}, e soprattutto
\textbf{perché}. La slide propone una tabella per generi
(tabella~\ref{gd18:tab:generi}), basata su una prima sfumatura del modello.

\begin{table}[ht]
\centering
\small
\begin{tabular}{@{}p{0.2\textwidth}p{0.42\textwidth}p{0.3\textwidth}@{}}
\toprule
\textbf{Genere} & \textbf{Giochi tipici} & \textbf{Stili di gioco principali} \\
\midrule
FPS & \emph{Halo}, \emph{Call of Duty}, \emph{Half-Life}, \emph{Crysis} & Killer,
Achiever \\
CRPG & \emph{Darklands}, \emph{Fallout 1/2}, \emph{Baldur's Gate} & Achiever,
Explorer \\
FPS-CRPG & \emph{Deus Ex}, \emph{BioShock}, \emph{Mass Effect} & Achiever,
Explorer, Killer \\
CRPG open world & \emph{The Elder Scrolls}, \emph{Fallout 3}, \emph{Two Worlds} &
Achiever, Explorer \\
MMORPG & \emph{World of Warcraft}, \emph{EVE Online}, \emph{Guild Wars} &
Achiever, Explorer, Socializer \\
MMOG & \emph{Unreal Tournament}, \emph{Team Fortress}, qualsiasi FPS multiplayer &
Killer, Achiever \\
Avventura & \emph{King's Quest}, \emph{Myst}, \emph{The Longest Journey} &
Socializer, Explorer \\
Azione & \emph{Tomb Raider}, \emph{Uncharted}, \emph{Angry Birds} & Killer,
Achiever \\
Survival horror & \emph{Resident Evil}, \emph{Dead Space}, \emph{Amnesia} & Killer,
Achiever \\
Strategia a turni & \emph{Civilization}, \emph{Master of Orion},
\emph{Galactic Civilizations} & Explorer \\
\bottomrule
\end{tabular}
\caption{Il modello prismatico applicato all'analisi dei giochi per genere (dalle
slide).}
\label{gd18:tab:generi}
\end{table}

\subsection{Il commento alla tabella}

\paragraph{FPS.} \emph{Call of Duty}, \emph{Half-Life}, \emph{Crysis}:
killer-achiever, ``sono d'accordo'', più o meno su tutti. Una tabella così è utile
per avere, in pre-produzione, un benchmark con una grande mole di dati da cui
partire. Ma se il competitor diretto è \emph{Half-Life}, serve sapere in modo più
chiaro che cosa lo distingue dagli altri: in \emph{Half-Life} c'è una componente
\textbf{explorer} importante, che \emph{Call of Duty} non ha (o ha molto meno).
Capire \textbf{quanta percentuale} di ciascun tipo c'è in un gioco, e tradurre questi
giudizi in \textbf{dati quantitativi}, è un passo che si farà più avanti. Per ora
basta individuare i \textbf{primi due player type dominanti}.

\paragraph{RPG.} Generalmente achiever-explorer, con l'explorer in testa; ma in un
gioco come \emph{RuneScape}\footnote{Nella trascrizione: ``BanduScape''; la lettura
\emph{RuneScape} è probabile.} c'è una componente socializer importante. In
\emph{Deus Ex}, \emph{BioShock} e \emph{Mass Effect} si trovano achiever, explorer e
killer; in \emph{BioShock}, poi, la relazione con le Sorelline (salvarle o
``massacrarle'') ha anche una componente socializer, ma è un elemento dentro un
gioco intero. Per \emph{The Elder Scrolls} (\emph{Skyrim}, \emph{Morrowind}) e
\emph{Fallout}, achiever-explorer, ``ma explorer su tutti'', e forse anche una
spiccata componente socializer.

\paragraph{MMO.} \emph{World of Warcraft}, \emph{EVE Online}, \emph{Guild Wars}:
achiever, explorer, socializer. \emph{EVE Online} ha un explorer molto spiccato,
mentre \emph{World of Warcraft} ``lo si gioca per stare con l'altra persona''.

\paragraph{La leggenda della tripla A.} Una leggenda, ``non so chi l'ha messa in
circolo, però siccome vive \emph{rent-free} nella mia testa, ve la contamino'': il
numero di A di un gioco non dipenderebbe solo dal costo di produzione, ma anche dal
\textbf{numero di player type} del suo audience. Un doppia A ne avrebbe due, un
tripla A tre; ora ci sono anche i quadrupla A, e \emph{GTA 6} li avrebbe tutti e
quattro. Ma non è vero nemmeno il contrario: anche \emph{The Binding of Isaac} ha
dentro tutti e quattro i player type, ``io sono sicuro'', solo che la proporzione di
socializer (e di qualche altro tipo) è estremamente ridotta; sui dubbi come questo
entra in gioco un'analisi più fine, per cui gli studenti non hanno ancora gli
strumenti, e anche il docente, per esserne sicuro, dovrebbe compilare un documento.
Per giochi molto piccoli, come \emph{Pong}, la componente degli altri tipi è così
ridotta che conviene dire solo il principale: \emph{Pong} ``sicuro non ha
l'explorer'', forse è achiever, ma se lo si gioca con un amico\dots

\paragraph{I giochi molto grandi.} In altri giochi le componenti sono tante, ma in
sistemi diversi. È ciò che si faceva in Ubisoft, nei giochi open world come
\emph{Assassin's Creed}: due player type principali, più una serie di gameplay
laterali, non obbligatori (missioni secondarie, \emph{side quest}), che
``stuzzicano'' l'interesse di altri tipi. Alla fine ci sono tutti, ma per i giochi
molto grandi interessa il \textbf{gameplay principale}: una sezione di pesca in uno
sparatutto non è da killer, ma è una componente molto contenuta rispetto a un monte
ore molto elevato. Dipende tutto da che cosa si considera; per semplicità si prendono
i \textbf{primi due}, individuando il principale e il secondario. Setacciare
precisamente le percentuali verrà più avanti.

\paragraph{Strategia a turni.} Non c'è solo l'explorer: a seconda dei titoli c'è
anche l'achiever. \emph{Slay the Spire} è sicuramente achiever; ma ha anche una
componente aleatoria molto elevata, che è da socializer, e le sfide, le restrizioni,
i livelli di ascensione più alti, che potrebbero interessare un killer. ``Dipende.''

% ---------------------------------------------------------------------
\section{Le feature di gameplay per player type}
\label{gd18:sec:feature}

Lo stesso ragionamento si può fare per \textbf{progettare giochi nuovi}: ci sono
feature di gameplay spiccatamente pensate per un player type, o spiccatamente vicine
a esso (tabella~\ref{gd18:tab:feature}). È un discorso che verrà espanso, sempre
legato alle percentuali.

\begin{table}[ht]
\centering
\small
\begin{tabular}{@{}p{0.3\textwidth}p{0.62\textwidth}@{}}
\toprule
\textbf{Stile di gioco unificato} & \textbf{Feature di gameplay associate} \\
\midrule
Artisan/Killer/Experientialist & azione, vertigine, uso di strumenti, uso di
veicoli, horror, gioco d'azzardo, speedrun, exploit \\
Guardian/Achiever/Gamist & competizione, collezioni, produzione, high score,
livelli, obiettivi chiari, appartenenza a una gilda, min-maxing \\
Rational/Explorer/Simulationist & puzzle, costruzione creativa, lore del mondo,
analisi dei sistemi, teorizzare, sorpresa \\
Idealist/Socializer/Narrativist & chat, roleplay, storytelling, cooperazione,
decorazione, animali domestici, eventi sociali \\
\bottomrule
\end{tabular}
\caption{Il modello prismatico applicato al design di giochi nuovi (dalle slide).}
\label{gd18:tab:feature}
\end{table}

\paragraph{Killer.} Usare veicoli, armi, esplosioni; l'horror; il \textbf{gamble}:
``rischio tutto, punto tutto sul 17'' è da killer, mentre tirare un dado per vedere
se si riesce a volare è da socializer (e il volare è da explorer). Le
\textbf{speedrun} sono da killer, come l'azione. Il docente prova anche a
categorizzare le motivazioni di Quantic Foundry: la distruzione è da killer, la
community da socializer, la completion da achiever; il \textbf{power} dipende da
come lo si ottiene: se tramite \emph{hard work}, achiever; se per sbilanciare il
gioco, killer. Fare l'impostore in \emph{Among Us} è da killer e da socializer.
L'immersione e la fantasy, se regolate, sono da explorer; la storia da socializer;
la discovery da explorer.

\paragraph{Self-expression.} Una cosa peculiare: la \textbf{customization}. Che cosa
potrebbe essere? Il poter esprimere sé stessi: mettere fiori e piante nella propria
casetta, farla verde, marrone o gialla, avere i capelli del colore o della forma che
si vuole. È molto vicina al socializer.

\paragraph{Achiever.} Tutto ciò che si \textbf{completa}: se c'è una barra, se c'è una
lista di cose da fare. I collezionabili; l'high score (che può valere anche per il
killer); i livelli, livellare, l'XP; gli obiettivi chiari; l'appartenenza a una
gilda; il \textbf{min-maxing}, che però è a cavallo: dipende se serve fare data
mining e che cosa si usa per farlo, e allora può essere anche explorer.

\paragraph{Explorer.} I \textbf{puzzle} su tutti; la lore, unire i puntini; la ricerca
dei dettagli, delle prove in un caso di omicidio; il \textbf{teorizzare}. E la
\textbf{sorpresa}, che però è ``un refuso'' del vecchio modello unificato da cui si è
partiti.

\begin{remark}[La sorpresa non è un'emozione]
La \textbf{sorpresa} è un \textbf{moltiplicatore emotivo}, non un'emozione. Se si è
sorpresi mentre si scappa dallo xenomorfo, con lo \emph{stalker} alle spalle, si è
ancora più spaventati: è il motivo per cui i \emph{jumpscare} a volte funzionano e a
volte no, a seconda che li si anticipi o meno. E se si spera di trovare in un
\emph{gacha} il personaggio ultra raro, con una probabilità molto bassa, e lo si
trova, si è estremamente contenti, perché la sorpresa potenzia l'anticipazione che
si aveva.
\end{remark}

Rientra nell'explorer anche esplorare non solo l'open world, ma il \textbf{come
funzionano} gli oggetti, lo sperimentare con le cose. \emph{Minecraft}, non a caso,
è da explorer: è un sandbox a disposizione, in cui vedere come le cose si combinano
(il \emph{crafting table}), come funzionano gli oggetti, come si muovono gli
ecosistemi. Non è l'unico; e \emph{Minecraft} cambia notevolmente player type se lo
si gioca in single player o in multiplayer, ma il primo player type resta l'explorer.

\paragraph{Socializer.} Tutto ciò che permette di \textbf{parlare} con un'altra
persona, di interagire in maniera esplicita: meglio se si parla, peggio se c'è solo la
chat scritta, ancora peggio se ci sono solo i \emph{ping}. Con opzioni limitate, un
ping finisce per diventare un insulto: il docente mostra l'emote più ``non
sportiva'' di \emph{Clash Royale} (ormai ne esistono versioni potenziate) e quella di
\emph{Marvel Snap} con Ms.\ Marvel, ``la cosa più innocente dell'universo'',
diventata l'insulto del gioco. Avere opzioni limitate crea accezioni specifiche per
quelle che ci sono; se ne parlerà in una lezione apposta. Poi lo storytelling, la
cooperazione, la decorazione e tutto ciò che è customizzabile, il \textbf{pet}, l'NPC
o il cagnolino che segue il giocatore.

Sono sistemi da cui partire: quelli di base, spiccatamente ben percepiti da uno o più
player type. ``Unite il tutto, create un vostro documento'': ce ne sono, e ce ne
saranno, molti altri da integrare nel modello.

\slidefig{img/GD/gd18-esprimere}{Esprimere i player type del target audience: elementi attraenti ed estetiche target per tipo.}{slide \emph{DM06 -- Persona \& Types}}{gd18:fig:esprimere}

Le slide mostrano anche un esempio di come \textbf{esprimere} i player type di un
target audience in un documento (figura~\ref{gd18:fig:esprimere}): per ogni tipo
(explorer, killer, achiever, con i colori del modello) gli \emph{appealing elements},
presi dalle motivazioni di Quantic Foundry e da Lazzaro (per esempio \emph{Fantasy
Deep Lore}, \emph{Discovery curious}, \emph{Easy Fun} per l'explorer;
\emph{Excitement fast-paced} per il killer; \emph{Power progression} e
\emph{Completion task-oriented} per l'achiever), e le \emph{target aesthetics}
corrispondenti (mondo alternativo, sandbox e open world, esplorazione; scarica di
adrenalina; partire deboli e salire di livello, stelle, trofei e collezionabili).

% ---------------------------------------------------------------------
\section{I segmenti di Quantic Foundry}
\label{gd18:sec:segmenti}

\slidefig{img/GD/gd18-segmenti}{I nove segmenti di giocatori di Quantic Foundry.}{slide \emph{DM06 -- Persona \& Types}}{gd18:fig:segmenti}

\subsection{Nove, non otto né dodici}

Il modello prismatico non è l'unico. Quantic Foundry, per esempio, ha un suo modello
di \textbf{nove} player type (figura~\ref{gd18:fig:segmenti}): ``nove, non otto né
dodici. Peccato''. Non sono ripartibili equamente sui quattro tipi, e questo, ammette
il docente, è un suo ``problema mentale'': quando deve categorizzare qualcosa, deve
essere diviso esattamente in parti uguali tra i player type; per lui un modello così
va integrato con altro, ma per gli studenti potrebbe non valere.

\begin{table}[ht]
\centering
\small
\begin{tabular}{@{}p{0.14\textwidth}p{0.26\textwidth}p{0.24\textwidth}p{0.27\textwidth}@{}}
\toprule
\textbf{Segmento} & \textbf{Motto} & \textbf{Motivazioni principali} &
\textbf{Giochi popolari} \\
\midrule
Acrobat & \emph{Flexing my reflexes} & challenge + discovery & \emph{Spelunky},
\emph{Celeste}, \emph{Super Metroid}, \emph{Tetris} \\
Gardener & \emph{Quiet, relaxing task completion} & completion & \emph{Candy
Crush}, \emph{Solitaire}, \emph{Animal Crossing} \\
Slayer & \emph{Cinematic mayhem with a purpose} & fantasy + story + destruction &
\emph{Firewatch}, \emph{Uncharted}, \emph{Tomb Raider} \\
Skirmisher & \emph{Jumping into the fray of battle} & destruction + competition &
\emph{Rust}, \emph{Call of Duty}, \emph{Battlefield} \\
Gladiator & \emph{Dedicated, hardcore gaming} & challenge + completion + community &
\emph{Mobile Legends}, \emph{Destiny}, \emph{Gears of War} \\
Ninja & \emph{A duel of speed and skill} & competition + challenge & \emph{Street
Fighter}, \emph{StarCraft}, \emph{LoL} \\
Bounty Hunter & \emph{High-octane solo world exploration} & destruction + fantasy &
\emph{Mass Effect}, \emph{Far Cry}, \emph{Saints Row} \\
Architect & \emph{My empire begins with this village} & strategy + completion &
\emph{Europa Universalis}, \emph{Civ VI}, \emph{Banished} \\
Bard & \emph{Playing a part in a grand story} & design + community + fantasy &
\emph{The Secret World}, \emph{FFXIV}, \emph{LoTRO} \\
\bottomrule
\end{tabular}
\caption{I segmenti di giocatori di Quantic Foundry (dalle slide).}
\label{gd18:tab:segmenti}
\end{table}

\subsection{Leggere i segmenti con Bartle}

\paragraph{Acrobat.} Giocatori solitari che cercano la sfida nel gameplay e si
esercitano ``\emph{over and over}'': già il \emph{practice over and over}, la sfida,
il single player dicono achiever. Challenge più discovery: achiever più explorer.

\paragraph{Il dibattito su Celeste.} Ma chi gioca \emph{Celeste}? Il docente chiede
agli studenti i due player type principali: ``non è facile''. Uno propone la sfida
del platform; un'altra fa notare che in \emph{Celeste} c'è una narrativa molto
profonda (``è vero, è vero''); un altro ancora parla di componenti sociali legate
alla sua mediaticità. Il docente è d'accordo sulla storia, ma ``non morirò su questa
collina'' (anche se poi, come sempre, bisogna trovare i dati): nel complesso
\emph{Celeste} è \textbf{achiever-killer}, perché ci sono una serie di sistemi
spiccatamente pensati per il killer: tutta la parte di \textbf{speedrun} e tutte le
sfide successive alla trama principale. La trama è explorer-socializer, c'è un
dramma importante; ma dopo c'è tutto il gameplay per cui \emph{Celeste} è ricordato,
e le sfide come \emph{Farewell}, il capitolo dopo la morte della nonna, ``tostissimo'':
quella roba è achiever-killer. ``Sto considerando tutto l'interesse del gioco.'' Se lo
si dà in mano al docente, finisce la campagna, fa qualche livello extra, ma non ce la
fa a fare \emph{Farewell}: ``sono quasi due giochi diversi''. Una studentessa racconta
di averlo platinato (``complimenti''); lui non ha le skill necessarie e non si sente
nemmeno spinto a coltivarle, perché ci sono tecniche come il \emph{wavedash} che sono
``toste''.

\paragraph{Metroid Dread.} È invece una cosa che il docente fa volentieri con
\emph{Metroid Dread}: i controlli sono ``atroci'', eppure l'ha completato e ha fatto
tutto. Il problema non è la modalità Dread, ma i puzzle dello \emph{shinespark}, che
richiedono troppe abilità; i comandi sono mappati non male, ma ``per necessità,
perché sono così tante'', e la mappatura ne risente: certe sfide sono terribili.

\paragraph{Spelunky e il Ninja.} Lo stesso vale per \emph{Spelunky}: qual è
l'obiettivo? Finire il gioco, battere la Sunken City, finirlo nel minor tempo
possibile? Dipende dall'obiettivo. Challenge più discovery è achiever-explorer;
la competizione, invece, è molto vicina al \textbf{Ninja}, che è killer-achiever. Un
giocatore di \emph{Spelunky} o di \emph{Celeste} il docente lo direbbe
killer-achiever, ``più Ninja che Acrobat''. Sono framework diversi, entrambi
ugualmente validi, e uno non esclude l'altro.

\paragraph{Gardener e Tom Nook.} \emph{Complete completion} è chiaramente achiever;
ma \emph{Animal Crossing}, che pure è nel Gardener, non è solo achiever. ``Quiet,
relaxing task completion'' è achiever-socializer, e forse ci sta come player type di
\emph{Animal Crossing}: dipende da che cosa vuole fare la gente nel gioco. Per il
docente, \emph{Animal Crossing} è all'origine della sua ansia: essere costantemente
in debito con \textbf{Tom Nook} lo ``disintegra''. Paghi il debito per la casa e ce
n'è un altro ancora più grande: ``mi manda fuori di testa. Vabbè, comunque sono io
malato''. Probabilmente, scherza, lo stato economico mondiale è messo meglio di quello
di \emph{Animal Crossing}.

\paragraph{Bounty Hunter.} Distruzione e fantasy: killer-explorer. \emph{Mass Effect}
è killer-explorer, con una componente socializer importante. E il gioco con il
rampino in cui si fa saltare tutto per aria? \emph{Just Cause}, suggeriscono gli
studenti: chiaramente killer-explorer.

\paragraph{Architect e Clash of Clans.} Strategia più completion: achiever-explorer,
se costruire un villaggio significa vedere come si relazionano le componenti. Ma
\emph{Clash of Clans} ha una componente sociale importante. Il primo player type è
sicuramente l'achiever; il secondo, per il docente, è il socializer: c'è la guerra tra
clan, ma nel loop di gioco si fanno meno guerre di quanto si interagisca nel clan.
Senza dubbio, però, anche andare a rubare le risorse agli altri fa parte del loop:
``anche quello, bisognerebbe vederlo''. Per ora gli strumenti bastano per una visione
``a volo d'uccello''; poi ne arriveranno altri. Gli altri segmenti il docente li
lascia agli studenti: ragionarci e cercare di incasellarli nei quattro tipi.

% ---------------------------------------------------------------------
\section{Come si usa il modello: tre casi}
\label{gd18:sec:casi}

``Ora, bello tutto. Come lo usiamo?'' Per il docente è ``già una figata'' così, ma gli
studenti potrebbero chiedersi come usarlo nel lavoro di tutti i giorni. Ci sono tre
casi principali (potrebbero essercene altri), che corrispondono a tre momenti del
ciclo iterativo di produzione: la definizione del problema, lo sviluppo e il test.

\slidefig{img/GD/gd18-caso1}{Primo caso: l'analisi dei competitor, all'inizio del ciclo iterativo.}{slide \emph{DM06 -- Persona \& Types}}{gd18:fig:caso1}

\subsection{Primo caso: l'analisi}

Il primo caso sta a monte della produzione (figura~\ref{gd18:fig:caso1}): capire il
problema del target audience, se vuole qualcosa e che cosa dargli di conseguenza. Una
volta fatta la ricerca, usando i player type si può capire che cosa piace a un tipo
di giocatore e che cosa no (anche al di là di ciò che non dice), che cosa si può
migliorare, come spende il tempo, quanto hanno venduto i giochi, soprattutto rispetto
ai \textbf{competitor}. Capire i player type principali di un gioco, o la
distribuzione dei giochi più venduti del mercato odierno rispetto ai player type, può
essere estremamente utile.

Si torna al casual di Bateman: allora il mercato era pieno di FPS, tutti
principalmente killer-achiever; ora è molto più vario, ma anche molto
\textbf{segmentato}. Capire per esempio che cosa vende e che cosa no nel mondo
indipendente (che è un gergo estremamente largo), o quali sono i player type
dominanti dei competitor principali o dei giochi più venduti, o di quelli che hanno
più notorietà, è utile. Usato come strumento di analisi, il modello aiuta a capire le
emozioni principali di un gioco, come sono divise e soprattutto come si distribuiscono
sui player type principali. Come farlo si vedrà più avanti.

\slidefig{img/GD/gd18-rimworld}{Un esempio di analisi: le emozioni e i player type di \emph{RimWorld}.}{slide \emph{DM06 -- Persona \& Types}}{gd18:fig:rimworld}

Al docente questa cosa piace molto, la usa quasi sempre: è molto dispendiosa, ma
porta validi risultati. Farla per un gioco completo, come \emph{RimWorld}
(figura~\ref{gd18:fig:rimworld}), che è estremamente grande, richiede molto tempo.
L'esempio delle slide mostra il risultato: un'analisi delle estetiche (con la
\emph{discovery}, la \emph{fantasy} come mondo credibile e la \emph{mimesis} come
emozioni principali, \emph{challenge} ed \emph{expression} come secondarie) e la
distribuzione dei player type, con l'\textbf{explorer} al 58\%, l'achiever al 24,5\%,
il socializer al 14,9\% e il killer al 2,6\%, ciascuno con le ragioni per cui il gioco
lo attrae. È il primo punto: l'analisi di giochi, di competitor, di meccaniche, di
sistemi.

\subsection{Secondo caso: i suggerimenti di design}

\slidefig{img/GD/gd18-caso2}{Secondo caso: suggerimenti di design a partire dalle parole chiave del modello.}{slide \emph{DM06 -- Persona \& Types}}{gd18:fig:caso2}

Il secondo caso è nello sviluppo. Si conoscono i player type della propria audience,
si ha la ricerca. Ora, tra l'altro, è chiaro a che cosa si riferiva il target
audience della vision: nella vision alcuni elementi saranno pensati spiccatamente per
alcuni player type, altri per altri. Mettiamo che si stia facendo un FPS,
principalmente per achiever: si ha la vision, si conoscono i player type dominanti, e
si sa che l'audience vuole qualcosa che ha a che fare con le
\textbf{biciclette}\footnote{``Potrebbe essere qualsiasi cosa'': la bicicletta è solo
un esempio volutamente arbitrario.}. Come si fa il design? Ed ecco la differenza con
il \emph{data driven} di cui si parlava all'inizio: come si \textbf{riduce la falce}
in modo da aumentare la sicurezza che ciò che si inserirà nel gioco piacerà al
target audience, mantenendo però un \textbf{margine creativo}?

Si va nella propria lista, si prendono le \textbf{parole chiave} e si cerca di
sviluppare meccaniche o sistemi a partire da ciò che dice la definizione di quel
player type (figura~\ref{gd18:fig:caso2}). Gli esempi (``sono idee molto terribili,
però'') sono tre:
\begin{description}
  \item[Guardian: \emph{hard working} e \emph{seek security}.] Ogni livello (o
        mappa) contiene un \textbf{pezzo di bicicletta}; quando il giocatore (o il
        team) li trova tutti, può usare biciclette speciali che lo rendono immune ai
        danni. Pezzi da raccogliere con fatica, che una volta uniti danno sicurezza.
  \item[Lazzaro: \emph{mastery} e \emph{challenge}.] L'abilità speciale del giocatore
        è \textbf{correre e sparare in bicicletta}. Permette di muoversi
        velocissimi, ma è difficile muoversi e sparare insieme (magari i comandi sono
        invertiti rispetto all'implementazione base di qualsiasi veicolo); in cambio,
        il danno critico è raddoppiato.
  \item[Achiever e \emph{role culture}: gerarchia e lungo termine.] Le azioni
        speciali danno punti, che si spendono per \textbf{potenziare la bicicletta}
        (``l'achiever classico''). La bicicletta è lo status, il \emph{rank}, del
        giocatore, e il rank permette di portarsi dietro un certo numero di NPC.
        Questo è legato alla \emph{role culture} di Handy, che però ``non è che
        dobbiate per forza sposare''.
\end{description}
La tabella a fianco, nella slide, raccoglie le parole chiave del guardian/achiever in
tutti i modelli: logistico, laborioso, leale, responsabile, prudente, fiducioso
nell'autorità, in cerca di sicurezza (Keirsey); accumulare tesori, raccogliere punti e
salire di livello, status, gerarchia, competizione (Bartle); competizione, disciplina,
perseveranza, regole applicate a tutti allo stesso modo (Caillois); mastery, sfida,
obiettivi, progresso (Lazzaro); e così via fino alla classe del \textbf{guerriero}, il
cui motto è ``\emph{I have not yet begun to fight!}''.

\subsection{Terzo caso: risolvere un problema di design}

\slidefig{img/GD/gd18-caso3-tabelle}{Terzo caso: le parole chiave di artisan/killer e di idealist/socializer, con quelle evidenziate.}{slide \emph{DM06 -- Persona \& Types}}{gd18:fig:caso3-tabelle}

L'ultimo caso è nella fase di iterazione, nel test. Si sa che qualcosa non funziona,
ma non si sa perché. L'esempio viene da uno studente del primo anno: un
\textbf{gioco di carte per killer e socializer}. I test mostrano che la
\textbf{fine della partita è noiosa}: non piace ai giocatori, è troppo simile al resto
del gioco ed è chiaro chi vincerà. Serve una feature per il finale che piaccia
\textbf{sia ai killer sia ai socializer}, perché sono i tipi di riferimento, e che
renda il finale \textbf{più incerto}.

Si prende la lista e si leggono le parole chiave (figura~\ref{gd18:fig:caso3-tabelle}).
Per i killer: \textbf{tattico}, \textbf{imposizione sugli altri} e \textbf{il brivido
della caccia} (\emph{thrill of the chase}), quindi adrenalina, rischio, azzardo. Per i
socializer: \textbf{diplomatico}, \textbf{roleplay}, \textbf{interessato alle
persone} e \textbf{la sorte} (\emph{chance}).\footnote{Il docente aggiunge che, in
retrospettiva, nelle slide evidenzierebbe altre parole, e lo farà quando le
aggiornerà.}

\paragraph{Un'interruzione.} A questo punto entra in aula un collega, coordinatore del
modulo di narrative design, che saluta la classe (``ci siamo già visti con molti di
voi''), annuncia le lezioni di narrativa che partiranno già il venerdì successivo e
augura ``buon player type'' prima di uscire.\footnote{La trascrizione di questo
passaggio è molto frammentaria: i nomi dei docenti di narrativa citati non sono
ricostruibili con certezza.}

\slidefig{img/GD/gd18-bluff}{Che cosa hanno in comune killer e socializer? Il bluff.}{slide \emph{DM06 -- Persona \& Types}}{gd18:fig:bluff}

\paragraph{Che cosa hanno in comune?} Si riprende: che cosa hanno in comune queste due
liste? La slide mostra un tavolo da poker. Gli studenti ci arrivano: puntare le fiches
può avere le sue ragioni, ma è più una strategia che una meccanica; e la
\emph{poker face}, perché la si fa? La risposta è il \textbf{bluff}
(figura~\ref{gd18:fig:bluff}). Nel poker c'è una strategia chiaramente pensata sia per
i socializer sia per i killer. Ai killer piace perché permette di \textbf{imporsi}, di
rischiare e di avere un picco di adrenalina; e anche mentire, parlare, comunicare fa
parte del gioco del socializer, oltre alla componente di sorte. ``Chiaramente, più
killer.''

\subsection{I tre usi, riassunti}

Si possono quindi usare i player type, a partire dalla propria lista (dal modello
unificato, da ciò che si vuole), in tre modi:
\begin{enumerate}
  \item nell'\textbf{analisi}: prendere una lista di giochi e identificarne i player
        type dominanti aiuta a capire l'andamento del mercato, che cosa piace ai
        giocatori o che cosa vende, quali sono i player type dominanti dei
        competitor principali (come farlo bene si capirà più avanti);
  \item nel \textbf{fare design}: unire la componente creativa con una componente
        \textbf{scientifica}, mantenendo aperta la possibilità di soluzioni molto
        creative (perché si ha un numero ``esagerato'' di definizioni da integrare),
        ma riducendo la possibilità che vengano percepite male dall'utente o non
        apprezzate;
  \item nel \textbf{test}, durante il processo iterativo, per risolvere i problemi.
\end{enumerate}
Il docente concede cinque minuti di pausa: poi si conclude il persona, ``che tanto è
semplice'', e le ultime informazioni sui player type.

% ---------------------------------------------------------------------
\section{Il persona finale}
\label{gd18:sec:persona-finale}

\slidefig{img/GD/gd18-persona-finale}{Il persona finale: i dettagli non sono inventati, ma rappresentano i tratti più comuni del target audience.}{slide \emph{DM06 -- Persona \& Types}}{gd18:fig:persona-finale}

\subsection{Quanto costa un persona}

Al rientro, il docente riprende il \textbf{persona}. Non se ne farà uno in aula, per
questioni di tempo; ma può essere utile arrivarci con la ricerca, e averne uno nella
documentazione può dare supporto. Richiede però \textbf{tanto sforzo}: non si può
pensare di scalfire solo la superficie, con una descrizione degli elementi
superficiali. E c'è di più: se si crea un persona, poi \textbf{bisogna seguirlo}.
Non lo si crea per poi dimenticarlo: è come se quell'individuo inesistente
\textbf{entrasse a far parte del team}, ``non come designer, ma come tester''. Come
tutte le cose, ha i suoi pro e i suoi contro.

\subsection{Non si inventa niente}

\begin{quote}
I dettagli del persona non sono inventati: rappresentano i tratti più comuni del
target audience. Se non trovate questi dati, dovete fare più ricerca. Potete fare
assunzioni, ma non potete inventare.
\end{quote}
Quanti soldi ha? Quanto tempo spende giocando? Qual è la sua routine? Ha figli?
Statisticamente parlando: se non si hanno questi dati, quella ricerca va fatta. Per
questo è difficile creare un persona sufficientemente profondo da avere senso nel
team senza impiegarci davvero tante risorse. Si può chiedere al target audience che
cosa gli piace e che cosa no (intervistarlo, raccogliere dati, leggere i suoi post sui
social) e poi fare una \textbf{commistione} di queste cose; ma non si possono
inventare. ``È una ricerca molto lunga, non sottovalutatela.'' Se la si fa, poi
obbliga: a monte bisogna dedicarle tanto tempo e tanta ricerca, e a valle bisogna
seguirla, perché si è prodotto qualcosa e se ne vuole un risultato vantaggioso.

\subsection{Che cosa contiene}

Gli elementi del persona finale (figura~\ref{gd18:fig:persona-finale}) sono:
\begin{description}
  \item[Dati demografici.] La \textbf{foto} (oggi la si può generare facilmente, ed è
        meglio: così si evita di ``incappare in un Mario Rossi che già esiste''; il
        persona deve essere inventato, non una persona reale), il nome, l'età, il
        genere, lo stato relazionale, la nazione. Tutte le demografiche che danno con
        estrema precisione un'idea di chi sia.
  \item[Profilo da gamer.] Assumendo di fare un videogioco (il metodo si usa anche per
        i film): i giochi preferiti, la \textbf{combinazione di player type} del
        modello prismatico (achiever, killer, explorer, socializer), le motivazioni di
        gioco (per esempio da Quantic Foundry, ``ne vedremo altre''), i giochi e le
        meccaniche odiate, le abitudini di gioco (compra tutti i giochi o no, è
        completista o no, fa speedrun o no).
  \item[Profilo di cultura pop.] Le IP preferite (musica, film, anime), che cosa
        ascolta, che cosa guarda, che cosa odia. Il docente fa un esempio personale
        su un anime che non gli piace e su un eventuale seguito.\footnote{Nella
        trascrizione: ``non mi piace il Jiu Jitsu Kaiser, se dovessero fare un
        attiacodini 2''; i titoli non sono ricostruibili.}
  \item[Stato socio-economico.] Lavora o no? Lavorare prende tempo, e il passo
        successivo è definire la sua routine: quando gioca, dove, come. Ha tanti soldi
        o pochi? Magari ne ha tanti e spende poco, o pochi e spende tanto, magari
        compulsivamente; fa acquisti in-game o no, e che cosa compra? Gioca con amici,
        familiari, figli?
  \item[Dati personali.] Hobby, peculiarità, caratteristiche uniche.
  \item[Giornata tipo.] La sua routine va descritta.
  \item[Desideri e obiettivi.] Che cosa vuole dal prodotto, e che cosa vuole dalla
        sua vita.
\end{description}
Di solito si aggiunge anche un \textbf{motto}: ``non è che siamo NPC'', ma una frase,
una citazione con cui risuona particolarmente.

Tutto questo va in un documento di solito di \textbf{due pagine}, leggibile,
stampabile, molto contenuto: proprio perché va stampato e appeso, il team lo guarda,
``è come se ci fosse, dovrebbe esserci sempre all'interno del team''. È molto simile
alla vision, e la \textbf{complementa}.

\subsection{Uno o più persona?}

Una studentessa chiede: il persona è quindi la \textbf{media} del target, basata su
dati di campionamento? ``Sì, una campionatura statistica.'' E se ne fanno anche altri,
più piccoli? Non è detto che se ne faccia uno solo (il docente lo aveva detto
all'inizio): a maggior ragione quando si ha una mole enorme di dati, non conviene
farne uno, ma \textbf{più di uno}. Però \textbf{non più di tre}: ``già tre sono
tante''. Di solito il risultato è \textbf{uno principale e due secondari}, oppure
\textbf{due complementari}. Ma è molto lavoro: non va sottovalutato.

% ---------------------------------------------------------------------
\section{La vision, rivista}
\label{gd18:sec:vision}

\slidefig{img/GD/gd18-vision}{La vision: chi, quando, dove (ricerca); che cosa e perché (scienza e arte); come (conoscenza).}{slide \emph{DM06 -- Persona \& Types}}{gd18:fig:vision}

Il docente torna sulla \textbf{vision} (capitolo~\ref{gd9:cap:agency-visione}) e la
rilegge alla luce della lezione (figura~\ref{gd18:fig:vision}).

\paragraph{Interamente basato sulla ricerca.}
\begin{itemize}
  \item \textbf{Chi}: il target audience. Il gioco è per il target audience,
        \textbf{non per voi}. Si potrebbe farlo per sé stessi, ma è un rischio, ``un
        rischio molto grande'', e non è così che si sviluppano prodotti multimediali
        con grossi team alle spalle.
  \item \textbf{Quando}: il time frame, le milestone, il tempo di produzione. La
        ricerca, rispetto ai soldi disponibili, agli investimenti, a ciò che si sa
        fare e a ciò che si ha a disposizione nel team, determina il time frame; a
        volte lo impone lo stakeholder (``fai il gioco in tot tempo'').
  \item \textbf{Dove}: le piattaforme, il team, la comunicazione, gli stakeholder.
\end{itemize}

\paragraph{50\% scienza + 50\% arte = game design.}
\begin{itemize}
  \item \textbf{Che cosa}: il concept, le emozioni, le meccaniche. Si parte dalla
        ricerca, poi si integrano i player type e qualsiasi framework si voglia usare
        come designer: questo incrocio dà il gioco.
  \item \textbf{Perché}: i \emph{design value}, gli obiettivi come sviluppatore. Si
        potrebbe voler fare tanti soldi; potrebbe anche essere basato completamente
        sulla ricerca. Un'azienda ha dei valori (non tutte, ma molte): il valore
        principale è il denaro, e va bene finché si sta a galla, il gioco deve
        vendere tanto, ed è un valore di per sé; ma per le aziende valide di solito
        non è l'unico.
\end{itemize}

\paragraph{Interamente basato sulla conoscenza.} Il \textbf{come}: ciò che si è
studiato (engine, agency, meccaniche, narrativa\dots), l'esperienza e il gusto. Che, a
modesto avviso del docente, intacca anche le altre parti: ciò che si sa fare abilita
alcune cose. Il che non significa che non si possa fare un gioco che non si sa fare:
il docente lavora costantemente a giochi che richiedono competenze che non ha del
tutto, o che deve imparare. L'importante è restare aperti di mente; e la difficoltà
del percorso in accademia aiuta anche in questo, perché superare molti ostacoli
facilita il domani. Dipende anche da come si è: si può voler essere più sicuri e
affrontare un titolo più contenuto, più vicino alle proprie competenze, aumentando la
probabilità che il gioco esca o di raggiungere l'obiettivo; oppure rischiare tanto,
imbarcarsi in un progetto che non si ha la certezza di saper fare e imparare durante
il percorso. Comunque vada, si impara. In accademia, soprattutto in analisi, il
rischio è zero; ``là fuori'' il rischio di consumare tempo e soprattutto soldi del
team è molto elevato.

\slidefig{img/GD/gd18-valori}{I design value, raggruppati per categoria.}{slide \emph{DM06 -- Persona \& Types}}{gd18:fig:valori}

\subsection{I design value}

Ci sono poi i \textbf{design value} (figura~\ref{gd18:fig:valori}): il docente non li
vede tutti, ma alcuni sono generalmente condivisi.
\begin{itemize}
  \item \textbf{Centralità del giocatore}: \emph{players' advocacy}, co-creatività e
        inclusione degli utenti, usabilità e giocabilità.
  \item \textbf{Game design casual}: accessibilità (facile da giocare e da
        acquistare), accettabilità, flessibilità, semplicità. Molti giochi mobile
        seguono questi valori, giustamente, perché ormai il casual gamer gioca
        principalmente lì.
  \item \textbf{Game design tradizionale}: immersione, sfida e competizione,
        comunità, \emph{other-worldliness}.
  \item \textbf{Espressione artistica, innovazione e sperimentazione}: visual design
        ed estetica, sperimentazione, \textbf{design divergente}. Quest'ultimo è ciò
        che si è visto nella lezione sul brainstorming
        (capitolo~\ref{cap:brainstorming-flowchart}) come limitazione dell'idea:
        prendere un'idea e stravolgerla completamente, lontano dal mercato attuale.
  \item \textbf{Valori ludologici}: divertimento e valore del fun, agnosticismo
        tecnologico, nostalgia ed estetica retrò, valore delle meccaniche di gioco.
  \item \textbf{Impatto sociale e valori culturali}: \emph{games for good} e giochi
        d'impatto, diversità e accessibilità, etica e moralità, diversità culturale e
        tradizione.
  \item \textbf{Processo di produzione e creazione}: rispetto tra pari e identità
        professionale, collaborazione e lavoro di squadra, ideologia open source, cura
        e dettagli, avanzamento tecnologico, lo sviluppo come sfida.
  \item \textbf{Indipendenza}: autonomia e libertà artistica, anarchia. Il valore
        dell'indipendenza non è necessariamente l'anarchia: è autonomia, completa
        libertà imprenditoriale, ``con i rischi del caso''.
  \item \textbf{Valori commerciali}: successo economico, opportunismo e
        \emph{disruption}.
\end{itemize}
L'obiettivo è averli chiari in testa e \textbf{comunicarli al team}. Durante il
progetto, legato alla vision, tutti devono condividere lo stesso obiettivo, e bisogna
\textbf{parlarne prima}: ``non è che lo si scopre poi che io volevo fare un sacco di
soldi e tu volevi sperimentare e imparare qualcos'altro, non funziona così''. Alcuni
valori sono conciliabili, altri meno; ma se non se ne parla prima, non ci si arriva.

% ---------------------------------------------------------------------
\section{La profilazione del target audience}
\label{gd18:sec:ta}

\slidefig{img/GD/gd18-ta-ricerca}{La profilazione del target audience: si parte dai dati sui competitor e sui giochi simili.}{slide \emph{DM06 -- Persona \& Types}}{gd18:fig:ta-ricerca}

La prima cosa da fare per un progetto, come già detto parlando di pre-produzione, è
\textbf{ottenere dati} (figura~\ref{gd18:fig:ta-ricerca}). Poi si guardano i
competitor, e non solo: siccome il mercato è molto ampio, ci si limita a giochi con
qualità simili a quelle del gioco che si vuole sviluppare:
\begin{itemize}
  \item \textbf{meccaniche} simili;
  \item \textbf{estetiche} simili;
  \item \textbf{pilastri} simili;
  \item \textbf{mercato} simile.
\end{itemize}
Si vuole vendere tanto? Si guardano i giochi più venduti. Si vuole vendere a un
mercato specifico? Si guarda che cosa piace a quel mercato. Si \textbf{riduce la
falce}, ma si fa tanta ricerca ad alto livello; e si possono considerare anche i
player type di questi giochi, identificandone il dominante.

\slidefig{img/GD/gd18-ta-fine}{Che cosa si deve avere alla fine della pre-produzione.}{slide \emph{DM06 -- Persona \& Types}}{gd18:fig:ta-fine}

\subsection{Alla fine della pre-produzione}

Alla fine della pre-produzione si deve avere (figura~\ref{gd18:fig:ta-fine}):
\begin{itemize}
  \item un \textbf{profilo specifico del target audience}, con le
        \textbf{motivazioni} e le \textbf{percentuali dei player type}, e le sue
        eventuali segmentazioni o sottocategorie;
  \item lo \textbf{stato del proprio segmento di mercato} (a seconda del genere, dei
        tag e simili): si vuole fare un metroidvania? Come sta il mercato dei
        metroidvania, c'è spazio o no? Generalmente lo si fa con un
        \textbf{grafico cartesiano} con gli assi, cercando gli \textbf{oceani blu}:
        combinazioni che non sono ancora state pensate, ammesso che non ci sia un
        motivo valido per cui non sono state fatte;
  \item il mercato di \textbf{distribuzione}: dove verrà venduto il gioco;
  \item i \textbf{dati} a supporto di tutto ciò che si afferma.
\end{itemize}

Sono i dati che serviranno nella lezione successiva, quella sul \textbf{pitching}.
Questa parte conclude le lezioni sulla pre-produzione: ci saranno altre nozioni sparse
qua e là, ma già ora gli studenti dovrebbero avere tutte le informazioni necessarie
per partire con la vision di un gioco e cominciare a mettere a terra un'idea a partire
dalla ricerca. Non tutta: pulita, ridotta, ``un pezzettino, dati chiari''. Questi
possono essere messi in un \textbf{pitch} per cominciare a vendere, cioè a chiedere o
a vendere qualcosa partendo dai dati. Alle spalle si ha una mole di dati, ``un drive
pieno di ricerche'', di documentazione, di grafici; e \textbf{i grafici sono il
risultato}, la cristallizzazione semplificata e molto visibile di tutto il lavoro che
sta dietro, fondamentale per proiettarsi verso il futuro.

% ---------------------------------------------------------------------
\section{Il calendario e la chiusura}
\label{gd18:sec:chiusura}

È la cosa che il docente aveva chiesto di ricordargli, e la ``smacchia subito''. La
prossima lezione con lui è molto lontana:
\begin{itemize}
  \item il \textbf{4} (glielo ha ricordato Daniele) c'è, da calendario,
        \textbf{narrativa}; per l'\textbf{11} il docente non ricorda, ``penso
        Alessandro'';
  \item il \textbf{14} c'è la lezione di \textbf{pitching} con \textbf{Diego e
        Jack}. Il docente non ci sarà; è una lezione molto importante, in cui
        confluiscono molte cose, soprattutto quella detta alla fine;
  \item durante la lezione del 14 verranno dati dei \textbf{task di pitching a
        gruppi}. ``So che è poco tempo'': il pitch va prodotto per il \textbf{18},
        quando si rivedranno con il docente e lo si presenterà in classe, per
        discuterne insieme. Non è riuscito a spostarlo al 21 (``non me la concede
        nessuno l'orario del 21''); è molto importante imparare a farlo prima di
        dicembre.
\end{itemize}
Venerdì prossimo, il 4, il docente sarà comunque in sede per la lezione al secondo
anno: la mattina, chi ha dubbi o domande può trovarlo.

Il ``lavoro'' resta quello detto: cominciare a lavorare sui player type, impararli,
familiarizzare con loro e con tutti i termini visti, e come cartina tornasole prendere
un po' di giochi che si conoscono bene e cercare di individuarne i principali player
type, e \textbf{perché}. Il docente caricherà nel materiale per gli studenti il
modello unificato, ``quello vecchio. Molto brutto apposta, quindi dovete lavorarci''.
Nessuna domanda: saluta Federica e ``i posteri che rivedranno questa lezione''.
``Ci vediamo il 18. Ciao.''

% ---------------------------------------------------------------------
\section*{Punti chiave del capitolo}
\addcontentsline{toc}{section}{Punti chiave del capitolo}

\begin{itemize}
  \item Non ci sono più intrusioni durante la lezione. Non si risolvono problemi se
        non si comunica: è il mantra del game design.
  \item La piramide: agency (limiti, non libertà), decisioni con outcome comunicato,
        feedback e conseguenze a breve termine, conseguenze a lungo termine,
        \textbf{mastery} dei sistemi. Attraverso l'esperienza il giocatore evolve
        verso il giocatore ideale del designer (i Gargoyle di \emph{Dark Souls}).
  \item Il gioco si fa \textbf{per qualcun altro}. Tre strumenti per capire
        l'audience: statistica, \textbf{metodo persona}, \textbf{player type}. Il
        design è \emph{data informed}, non \emph{data driven} (nel 90\% dei casi); lo
        skinning si fa, ma lo unique selling point deve stare altrove.
  \item Maslow: i videogiochi toccano i tre gradini più alti (autorealizzazione,
        appartenenza, stima), più l'esperienza vicaria. Il modello non spiega tutto
        (le notti in bianco, i bisogni ``vitaminici''). Cibo, sesso e violenza
        catturano l'attenzione, ma l'effetto cala se usati troppo.
  \item Le motivazioni di Quantic Foundry variano per cultura (la Cina più challenge,
        gli Stati Uniti più completion): i dati sono lontani tra loro e serve
        un'approssimazione.
  \item Il persona è un utente \textbf{fittizio}, basato sul target audience, non su
        una persona reale né sul designer; i dettagli sono tutto.
  \item Bartle (da sapere a memoria): \textbf{killer} (rosso, acting on players),
        \textbf{achiever} (blu, acting on world), \textbf{explorer} (verde,
        interacting with world), \textbf{socializer} (giallo, interacting with
        players). Nessuno è un solo tipo; conta \textbf{perché} si usa una meccanica,
        non la meccanica.
  \item Intento esplicito e implicito: politician/griefer, planner/opportunist,
        scientist/hacker, networker/friend. Single player e multiplayer: killer e
        achiever si incontrano sulla leaderboard (ma con mezzi diversi), explorer e
        socializer sulla storia (lore regolata contro dramma).
  \item Keirsey: artisan (sensazioni, tattico), guardian (sicurezza, logistico),
        rational (conoscenza, strategico), idealist (identità, diplomatico); le
        sedici personalità si dividono per quattro. Bateman aggiunge i tipi
        intermedi (conqueror, manager, wanderer, participant, hardcore, casual), ma
        il suo casual non è quello di oggi.
  \item Caillois: \textbf{agon} achiever, \textbf{alea} socializer (rende tutti
        pari), \textbf{mimicry} explorer (la maschera ha regole), \textbf{ilinx}
        killer. Lazzaro (più importante di Caillois): serious fun killer, hard fun
        achiever (il più vicino al fiero), easy fun explorer, people fun socializer.
  \item Il modello prismatico allinea questi e altri modelli (Handy, Gallup, Covey,
        motivazione, problem solving, obiettivo: fare, avere, conoscere, diventare).
        Niente MDA+ (è l'MDA che si appoggia al modello, non il contrario) e
        sconsigliato GNS+.
  \item Il ``non task'': rifarsi il proprio modello unificato e usarlo in ogni task;
        identificare i due player type principali dei propri giochi preferiti, e
        perché.
  \item Per genere e per feature: la sorpresa è un moltiplicatore emotivo, non
        un'emozione; nei grandi giochi conta il gameplay principale; i segmenti di
        Quantic Foundry si leggono anch'essi con Bartle (\emph{Celeste}
        achiever-killer per le sfide post-game).
  \item Tre usi del modello: \textbf{analisi} (competitor, mercato; \emph{RimWorld}),
        \textbf{suggerimenti di design} (le biciclette), \textbf{soluzione di
        problemi} (il finale del gioco di carte: il bluff).
  \item Il persona finale: dati non inventati (assunzioni sì, invenzioni no), foto
        generata, demografia, profilo da gamer, cultura pop, stato socio-economico,
        giornata tipo, desideri, motto; due pagine stampabili; al massimo tre persona.
  \item La vision: chi, quando, dove (ricerca); che cosa e perché (50\% scienza,
        50\% arte); come (conoscenza). I design value vanno condivisi con il team
        prima di cominciare.
  \item Alla fine della pre-produzione: profilo del TA con percentuali dei player
        type, stato del segmento di mercato (oceani blu), distribuzione, dati a
        supporto. Il 14 pitching con Diego e Jack, il 18 si presenta il pitch in
        classe.
\end{itemize}
